\chapter{Calcolo deduttivo proposizionale}

Nel capitolo precedente abbiamo definito una logica come una tripla \((\mathcal{L}, \models, \vdash)\), dove \(\mathcal{L}\) è un linguaggio, \(\models\) è la relazione di soddisfacibilità che ne esprime la semantica, e \(\vdash\) è un calcolo, che ci permette di dedurre formule a partire da altre, o verificare se una formula è valida o meno.

Per quanto concerne la logica proposizionale, abbiamo definito il linguaggio \(\PL\) e fornito ben due semantiche, ma nonostante abbiamo visto una forma di calcolo tramite le tabelle di verità, questo non è sufficiente.
Infatti, questo tipo di calcolo non \qi{scala} a logiche più complesse ed espressive, poiché assume il potersi ridurre a un numero finito di casi, cosa che già per il prim'ordine non è possibile.

Vedremo in questo capitolo due calcoli, un primo detto \keyword{calcolo di deduzione naturale} più affine al ragionamento umano, e un secondo detto \keyword{calcolo di risoluzione} più affine a un'implementazione automatica, e che è alla base di molti sistemi di dimostrazione automatica.

Entrambi questi calcoli sono di tipo puramente sintattico, e non fanno riferimento alla semantica, ma le regole di trasformazione che li caratterizzano sono formulate in modo da preservare le proprietà semantiche, e quindi ci permettono di dedurre conseguenze logiche, ovvero di ottenere la stessa relazione di deducibilità che abbiamo definito tramite la semantica, ma in maniera puramente sintattica.

\begin{definition}{Calcolo inferenziale}
  Definiamo un \keyword{calcolo inferenziale} come una coppia \((\Axi, \Rules)\) dove \(\Axi \subseteq \PL\) è un insieme di formule detto \keyword{assiomi}, e \(\Rules\) è un insieme di regole di trasformazione dette \keyword{regole di inferenza}.

  Le formule in \(\Axi\) sono proposizioni considerate vere a priori, e le regole in \(\Rules\) sono regole che ci permettono a partire da un numero finito di formule, dette \keyword{premesse}, di dedurre una nuova formula, detta \keyword{conclusione}.
\end{definition}

Ovviamente, come anticipato non tutte le regole di trasformazioni sono regole d'inferenza valide, poiché vorremmo che queste regole preservassero delle proprietà semantiche.

A prescindere da ciò però, a partire da un calcolo inferenziale \((\Axi, \Rules)\) possiamo definire una relazione di deducibilità.

\begin{definition}[namedlabel={def:deducibility}{Deducibilità}]{Deducibilità}
  Dati un calcolo inferenziale \((\Axi, \Rules)\) e un insieme di formule \(\Gamma\), diciamo che \(\Gamma\) deduce una formula \(\phi\) tramite \(\Rules\), e scriviamo \(\Gamma \vdash_{\Rules} \phi\), se è possibile dedurre come conclusione \(\phi\) a partire da un numero finito di formule in \(\Gamma\) e \(\Axi\), applicando un numero finito di regole in \(\Rules\).

  Per indicare l'uso di una specifica regola di inferenza \(\Rule \in \Rules\), scriveremo \(\Gamma \vdash_{\Rule} \phi\) per indicare che \(\Gamma\) deduce \(\phi\) tramite una singola applicazione di \(\Rule\).
  In caso non ci sia ambiguità sul calcolo inferenziale, scriveremo semplicemente \(\Gamma \vdash \phi\).
\end{definition}

Intuitivamente, l'obiettivo che vorremmo ottenere con un calcolo inferenziale, e quindi con la relazione di deducibilità, sarebbe quello di catturare, in maniera puramente sintattica, la conseguenza logica, ovvero vorremmo che:
\[\Gamma \models \phi \iff \Gamma \vdash \phi\]

Per quanto riguarda gli assiomi, sono un insieme minimale di tautologie dalle quali tramite il teorema di sostituzione varranno in diverse istanziazioni.\footnote{In un certo senso la regola di sostituzione fa parte di \(R\).}

Il primo calcolo che vedremo però non ha assiomi, ma solo regole di inferenza, e si chiama \keyword{calcolo di deduzione naturale}.

\section{Calcolo di deduzione naturale}

Avendo già anticipato l'assenza di assiomi, il calcolo di deduzione naturale esclusivamente dalle regole di inferenza, che sono formulate in modo da catturare il significato semantico degli operatori logici, e quindi ci permettono di dedurre conseguenze logiche.\footnote{
  La versione del calcolo di deduzione naturale presentato qui non è la sua versione minimale, e conterrà delle regole che non sono strettamente necessarie, ma che rendono più semplice la deduzione di formule, e quindi più affine al ragionamento umano.}

In particolare, per ogni operatore logico avremo una regola detta di \keyword{introduzione}, che ci permetterà di dedurre una formula che ha quell'operatore come connettivo principale, e una regola detta di \keyword{eliminazione}, che ci permetterà di dedurre formule che non hanno quell'operatore come connettivo principale, a partire da formule che invece lo hanno.

Nel nostro caso, considereremo i seguenti operatori logici: \(\land\), \(\lor\), \(\to\), \(\neg\) e \(\bot\), dove \(\bot\) è un operatore \(0\)-ario che rappresenta la contraddizione.

\begin{note}{}
  La notazione presentata è da intendersi come segue: sopra la linea ci sono le premesse, sotto la linea c'è la conclusione, e a destra della regola c'è il nome della regola, che è puramente indicativo, e serve solo a identificare la regola, ma non ha alcun significato formale.

  Inoltre, quando ci sono delle assunzioni che è possibile scaricare, queste sono indicate tra parentesi quadre, mentre i puntini di sospensione indicano l'applicazione di un numero finito di regole da quelle assunzioni fino alla conclusione.
\end{note}

\paragraph{Regole per \(\land\):}

Le regole di introduzione ed eliminazione per \(\land\) sono le più intuitive, e sono le seguenti:

\begin{equation*}
  \infer[(\land I)]{A \land B}{A & B} \quad\quad\quad \infer[(\land E^l)]{A}{A \land B} \quad\quad\quad \infer[(\land E^r)]{B}{A \land B}
\end{equation*}

\paragraph{Regole per \(\to\):}

Per quanto riguarda \(\to\), la regola di eliminazione è abbastanza intuitiva oltre che a essere molto conosciuta col nome di \textit{Modus Ponens}, mentre la regola di introduzione è un po' più complessa.

\begin{equation*}
  \infer[(\to E)]{B}{A \to B & A} \quad\quad\quad
  \infer[(\to I)]{A \to B}{
    \begin{array}{c}
    [A] \\
    \vdots \\
    B
  \end{array}
    } 
\end{equation*}

Quello che ci dice è che, se assumendo \(A\) riusciamo a dedurre \(B\), allora possiamo dedurre \(A \to B\) senza più assumere \(A\) \keyword{scaricando l'assunzione}. Tale regola cattura il procedimento usuale di dimostrazione di un teorema che prevede un'implicazione, ovvero volendo dimostrare \(A \to B\), allora si assume \(A\) e si tenta di dedurre \(B\). In caso affermativo, si potrà concludere \(A \to B\).

\paragraph{Regole per \(\lor\):}

Le regole di introduzione per \(\lor\) sono abbastanza intuitive, mentre la regola di eliminazione è un po' più complessa.

\begin{equation*}
  \infer[(\lor I^l)]{A \lor B}{A} \quad\quad\quad \infer[(\lor I^r)]{A \lor B}{B} \quad\quad\quad
  \infer[(\lor E)]{C}{
    \begin{array}[b]{c}
    A \lor B
    \end{array} &
    \begin{array}[b]{c}
    [A] \\
    \vdots \\
    C
  \end{array} &
    \begin{array}[b]{c}
    [B] \\
    \vdots \\
    C
  \end{array}
  }
\end{equation*}

Quest'ultima regola ci dice che, se abbiamo \(A \lor B\), e se assumendo \(A\) e \(B\) \textbf{separatamente} riusciamo a dedurre \(C\) in entrambi i casi, allora possiamo dedurre \(C\) senza più assumere \(A\) e \(B\), ovvero scaricando entrambe le assunzioni nelle rispettive deduzioni.
Questa regola cattura il concetto di dimostrazione per casi, ovvero se si vuole dimostrare \(C\) a partire da \(A \lor B\), sarà sufficiente dedurre \(C\) sia nel caso in cui sia vero \(A\), che nel caso in cui sia vero \(B\). In caso affermativo, si potrà concludere \(C\).
 
\paragraph{Regole per \(\neg\):}

Le regole di introduzione ed eliminazione per \(\neg\) sono le seguenti:

\begin{equation*}
  \infer[(\neg I)]{\neg A}{
    \begin{array}{c}
    [A] \\
    \vdots \\
    \bot
  \end{array}
  } \quad\quad\quad
  \infer[(\neg E)]{\bot}{\neg A & A}
\end{equation*}

La prima regola ci dice che, se assumendo \(A\) riusciamo a dedurre \(\bot\), allora possiamo dedurre \(\neg A\) senza più assumere \(A\) scaricando l'assunzione. Tale regola cattura il procedimento usuale di dimostrazione di un teorema che prevede una negazione, ovvero volendo dimostrare \(\neg A\), allora si assume \(A\) e si tenta di dedurre \(\bot\). In caso affermativo, si potrà concludere \(\neg A\).

\paragraph{Regola di riduzione all'assurdo:}

Infine, vediamo due regole che potremmo definire di eliminazione di \(\bot\), di cui una generalizzazione dell'altra.

\begin{equation*}
  \infer[(\bot E)]{A}{\bot} \quad\quad\quad
  \infer[(\RAA)]{A}{
    \begin{array}{c}
    [\neg A] \\
    \vdots \\
    \bot
  \end{array}
  }
\end{equation*}

La seconda regola è esattamente la \textit{Reductio ad Absurdum}, e ci dice che, se assumendo \(\neg A\) riusciamo a dedurre \(\bot\), allora possiamo dedurre \(A\) senza più assumere \(\neg A\) scaricando l'assunzione. Tale regola cattura il procedimento usuale di dimostrazione per assurdo, in cui negando la tesi si riesce a concludere una contraddizione. L'altra regola invece, che è una generalizzazione della seconda, cattura il concetto di implicazione da premesse false.
Essenzialmente quello che esprime è che, a partire da premesse contraddittorie, è possibile dedurre qualsiasi formula, e quindi in particolare \(A\). Un analogia per avere più chiara una giustificazione intuitiva di questa regola è considerare la semantica usuale dell'implicazione, nella quale, con antecedente falso, non importa il valore di verità del conseguente, ovvero a partire da premesse false è vero tutto e il contrario di tutto.

Il processo di deduzione naturale può quindi essere visto come una costruzione incrementale: si parte da un insieme di assunzioni iniziali e, mediante passi atomici costituiti da applicazioni di regole di inferenza, si ottiene progressivamente la formula obiettivo.
Durante tale processo, alcune assunzioni vengono scaricate tramite le regole che lo consentono (ad esempio \((\to I)\), \((\neg I)\), \((\lor E)\), \((\RAA)\)), mentre le altre restano attive fino alla fine della derivazione.
Queste ultime costituiscono le assunzioni residue, e sono precisamente quelle che compaiono a sinistra del simbolo di deducibilità.
In altri termini, una derivazione conclude \(\Gamma \vdash \phi\) quando \(\phi\) è ottenuta a partire da assunzioni non scaricate esattamente pari a \(\Gamma\).

Vediamo quindi una formalizzazione più rigorosa di questo processo, e in particolare di come si costruisce un albero di deduzione.

\begin{definition}{Albero di deduzione ben formato}
  Un \keyword{albero di deduzione} è un albero finito in cui nodi sono etichettati con formule di \(\PL\).
  
  Un tale albero è detto \keyword{ben formato} se ogni nodo interno è ottenuto dai suoi figli mediante un'istanza di una regola del calcolo naturale.
  
  Le etichette delle foglie sono \keyword{occorrenze di assunzioni}.
  Se un nodo interno è ottenuto da una regola che scarica un'assunzione, allora tutte le occorrenze di quell'assunzione nel sottoalbero radicato in quel nodo sono dette \keyword{scaricate}.
\end{definition}

\begin{definition}{Deduzione nel calcolo naturale}
  Dato \(\Gamma\subseteq \PL\) e \(\phi\in\PL\), il \keyword{sequente} \(\Gamma \vdash \phi\) è detto \keyword{deducibile} se esiste un albero di deduzione ben formato con radice etichettata da \(\phi\) e foglie etichettate da formule in \(\Gamma\) o assunzioni scaricate.
\end{definition}

Andiamo ora a vedere un primo esempio di deduzione naturale, per poi procedere con esempi più complessi che ci aiuteranno a costruire l'intuizione necessaria per affrontare deduzioni più complesse, guidando la scelta delle regole da applicare.

\begin{example}[namedlabel={ex:natural_deduction}{Deduzione Naturale}]{Deduzione naturale}
  Supponiamo di voler dimostrare che \(\vdash A \to \neg \neg A\).\footnote{Indichiamo con \(\vdash A\) per indicare che \(A\) è dimostrabile senza ipotesi, ovvero che \(A\) è un teorema.}

  Chiaramente \(A \equiv \neg \neg A\), quindi ci aspettiamo di poter dimostrare questa formula.

  Iniziamo assumendo \(A\) e \(\neg A\), quindi:
  \begin{prooftree}\label{prtree:natural_deduction_example}
    \AxiomC{\({[A]}_2\)}
    \AxiomC{\({[\neg A]}_1\)}
    \BinaryInfC{\(\bot\)}
    \RuleLabel{\(\neg I_1\)}
    \UnaryInfC{\(\neg\neg A\)}
    \RuleLabel{\(\to I_2\)}
    \UnaryInfC{\(A \to \neg\neg A\)}
  \end{prooftree}

  Andiamo a etichettare le inferenze scaricate con i numeri delle regole di inferenza che le scaricano, in modo da poterle identificare facilmente.

  Possiamo vedere la rappresentazione come albero della deduzione come segue:
  \begin{center}
    \begin{tikzpicture}
      \node (root) {\(A \to \neg\neg A\)};
      \node (n1) [above =of root] {\(\neg\neg A\)};
      \node (n2) [above =of n1] {\(\bot\)};
      \node (n3) [above left=of n2] {\([A]\)};
      \node (n4) [above right=of n2] {\([\neg A]\)};
      \draw (n1) -- (root);
      \draw (n2) -- (n1);
      \draw (n3) -- (n2);
      \draw (n4) -- (n2);
    \end{tikzpicture}
  \end{center}
\end{example}

\subsection{Applicazione del calcolo naturale proposizionale}

In questa sezione andremo ad analizzare nel dettaglio alcuni esempi di deduzioni naturali che esploreranno tutte le regole di inferenza che abbiamo introdotto, in modo da avere un'idea più chiara di come funziona questo calcolo, e di come scegliere le regole da applicare per dedurre formule più complesse.

\subsubsection*{Dimostrazione di \(\vdash A \leftrightarrow \neg \neg A\).}


Sebbene nel calcolo non abbiamo considerato regole di deduzione per \(\leftrightarrow\), possiamo considerare sempre la formulazione equivalente mediante \(\land\) e \(\to\), ovvero \((A \to \neg \neg A) \land (\neg \neg A \to A)\).

Possiamo facilmente notare che per dimostrare una formula con \(\land\) come connettivo principale, è ragionevole pensare di usare \(\land I\), che ci permette di dedurre una formula con \(\land\) come connettivo principale a partire da due formule che hanno come connettivo principale i due operandi di \(\land\). Quindi come primo passo, possiamo pensare di usare \(\land I\) per ottenere \((A \to \neg \neg A) \land (\neg \neg A \to A)\) a partire da \(A \to \neg \neg A\) e \(\neg \neg A \to A\).

\begin{prooftree}
  \AxiomC{\(A \to \neg \neg A\)}
  \AxiomC{\(\neg \neg A \to A\)}
  \RuleLabel{\(\land I\)}
  \BinaryInfC{\((A \to \neg \neg A) \land (\neg \neg A \to A)\)}
\end{prooftree}

Ovviamente, \(A \to \neg \neg A\) è già stato dimostrato in \Cref{prtree:natural_deduction_example}, quindi possiamo banalmente sostituire quella foglia con l'albero di deduzione che abbiamo già costruito, e quindi ci resta da dimostrare \(\neg \neg A \to A\).

\begin{prooftree}
  \AxiomC{\Namedcref{ex:natural_deduction}}
  \UnaryInfC{\(A \to \neg \neg A\)}
  \AxiomC{\(\neg \neg A \to A\)}
  \RuleLabel{\(\land I\)}
  \BinaryInfC{\((A \to \neg \neg A) \land (\neg \neg A \to A)\)}
\end{prooftree}

Dobbiamo quindi dimostrare \(\neg \neg A \to A\). Per fare questo, possiamo andare a retroso, e pensare di usare la regola \(\to I\), che ci permette di dedurre un'implicazione a partire dal suo conseguente. Inoltre tale regola è particolarmente importante perché ci \qi{regala} un'assunzione scaricabile, ovvero ci permette di usare liberamente come assunzione il suo antecedente nel sottoalbero radicato in essa, avendo la possibilità di scaricarla alla fine della deduzione. Le assunzioni scaricabili possono essere da guida per la scelta dei passi successivi, poiché se abbiamo un'assunzione scaricabile di una formula \(A\), allora è ragionevole pensare di usare tale formula come premessa per dedurre altre formule, e quindi di usare regole di inferenza che hanno \(A\) come premessa.

\begin{prooftree}
  \AxiomC{\Namedcref{ex:natural_deduction}}
  \UnaryInfC{\(A \to \neg \neg A\)}
  \AxiomC{\(A\)}
  \RuleLabel{\(\to I_1\)}[\({[\neg\neg A]}_1\)]
  \UnaryInfC{\(\neg \neg A \to A\)}
  \RuleLabel{\(\land I\)}
  \BinaryInfC{\((A \to \neg \neg A) \land (\neg \neg A \to A)\)}
\end{prooftree}

A questo punto dobbiamo dimostrare \(A\), sapendo di poter usare liberamente \(\neg \neg A\) come assunzione. Osservando le regole che abbiamo a disposizione, nessuna ci permette direttamente di dedurre \(A\) da \(\neg \neg A\), ma \(\RAA\) ci permette di dedurre \(A\) da \(\bot\) scaricando l'assunzione di \(\neg A\).

\begin{prooftree}
  \AxiomC{\Namedcref{ex:natural_deduction}}
  \UnaryInfC{\(A \to \neg \neg A\)}
  \AxiomC{\(\bot\)}
  \RuleLabel{\(\RAA_2\)}[\({[\neg A]}_2\)]
  \UnaryInfC{\(A\)}
  \RuleLabel{\(\to I_1\)}[\({[\neg\neg A]}_1\)]
  \UnaryInfC{\(\neg \neg A \to A\)}
  \RuleLabel{\(\land I\)}
  \BinaryInfC{\((A \to \neg \neg A) \land (\neg \neg A \to A)\)}
\end{prooftree}

A questo punto il gioco è fatto, poiché ci resta da dimostrare \(\bot\), ma abbiamo come assunzioni scaricabili sia \(\neg A\) che \(A\), e quindi possiamo trivialmente usare \(\neg E\).

\begin{prooftree}
  \AxiomC{\Namedcref{ex:natural_deduction}}
  \UnaryInfC{\(A \to \neg \neg A\)}
  \AxiomC{\({[\neg\neg A]}_1\)}
  \AxiomC{\({[\neg A]}_2\)}
  \RuleLabel{\(\neg E\)}
  \BinaryInfC{\(\bot\)}
  \RuleLabel{\(\RAA_2\)}
  \UnaryInfC{\(A\)}
  \RuleLabel{\(\to I_1\)}
  \UnaryInfC{\(\neg \neg A \to A\)}
  \RuleLabel{\(\land I\)}
  \BinaryInfC{\((A \to \neg \neg A) \land (\neg \neg A \to A)\)}
\end{prooftree}

\subsubsection*{Dimostrazione di \((P\to Q)\to((Q\to R) \to (P\to R))\).}

Anche qui, il connettivo principale è \(\to\), che vogliamo dimostrare partendo da un insieme vuoto di premesse, quindi dovremo scaricare ogni assunzione che facciamo.
Possiamo ragionevolmente iniziare con \(\to I\) per ottenere la nostra formula a partire da \((Q\to R) \to (P\to R)\) e scaricando l'assunzione di \(P\to Q\).

\begin{prooftree}
  \AxiomC{\((Q\to R) \to (P\to R)\)}
  \RuleLabel{\(\to I_1\)}[\({[P\to Q]}_1\)]
  \UnaryInfC{\((P\to Q)\to((Q\to R) \to (P\to R))\)}
\end{prooftree}

Ci troviamo adesso a dover dimostrare un altra implicazione, e quindi possiamo pensare di usare nuovamente \(\to I\) per ottenere \((Q\to R) \to (P\to R)\) a partire da \(P\to R\) scaricando l'assunzione di \(Q\to R\).

\begin{prooftree}
  \AxiomC{\(P\to R\)}
  \RuleLabel{\(\to I_2\)}[\({[Q\to R]}_2\)]
  \UnaryInfC{\((Q\to R) \to (P\to R)\)}
  \RuleLabel{\(\to I_1\)}[\({[P\to Q]}_1\)]
  \UnaryInfC{\((P\to Q)\to((Q\to R) \to (P\to R))\)}
\end{prooftree}

Ancora una volta, ci troviamo a dover dimostrare un'altra implicazione, e quindi possiamo pensare di usare nuovamente \(\to I\) per ottenere \(P\to R\) a partire da \(R\) scaricando l'assunzione di \(P\).

\begin{prooftree}
  \AxiomC{\(R\)}
  \RuleLabel{\(\to I_3\)}[\({[P]}_3\)]
  \UnaryInfC{\(P\to R\)}
  \RuleLabel{\(\to I_2\)}[\({[Q\to R]}_2\)]
  \UnaryInfC{\((Q\to R) \to (P\to R)\)}
  \RuleLabel{\(\to I_1\)}[\({[P\to Q]}_1\)]
  \UnaryInfC{\((P\to Q)\to((Q\to R) \to (P\to R))\)}
\end{prooftree}

Adesso, facendoci guidare dalle assunzioni scaricabili, possiamo notare che usando \(\to E\) possiamo dedurre \(R\) a partire da \(Q\to R\) assumendo \(Q\).

\begin{prooftree}
  \AxiomC{\(Q\)}
  \AxiomC{\({[Q\to R]}_2\)}
  \RuleLabel{\(\to E\)}
  \BinaryInfC{\(R\)}
  \RuleLabel{\(\to I_3\)}[\({[P]}_3\)]
  \UnaryInfC{\(P\to R\)}
  \RuleLabel{\(\to I_2\)}
  \UnaryInfC{\((Q\to R) \to (P\to R)\)}
  \RuleLabel{\(\to I_1\)}[\({[P\to Q]}_1\)]
  \UnaryInfC{\((P\to Q)\to((Q\to R) \to (P\to R))\)}
\end{prooftree}

Se terminassimo ora la dimostrazione, avremmo una ottenuto una deduzione per il sequente \(Q\vdash (P\to Q)\to((Q\to R) \to (P\to R))\), poiché l'unica assunzione non scaricata è \(Q\). Usando però le assunzioni scaricabili, possiamo dedurre \(Q\) usando ancora una volta \(\to E\) a partire da \(P\to Q\) e \(P\).

\begin{prooftree}
  \AxiomC{\({[P\to Q]}_1\)}
  \AxiomC{\({[P]}_3\)}
  \RuleLabel{\(\to E\)}
  \BinaryInfC{\(Q\)}
  \AxiomC{\({[Q\to R]}_2\)}
  \RuleLabel{\(\to E\)}
  \BinaryInfC{\(R\)}
  \RuleLabel{\(\to I_3\)}
  \UnaryInfC{\(P\to R\)}
  \RuleLabel{\(\to I_2\)}
  \UnaryInfC{\((Q\to R) \to (P\to R)\)}
  \RuleLabel{\(\to I_1\)}
  \UnaryInfC{\((P\to Q)\to((Q\to R) \to (P\to R))\)}
\end{prooftree}

Solo ora, possiamo concludere la dimostrazione, poiché ogni assunzione è stata scaricata.

\subsubsection*{Dimostrazione di \(\neg(P\land\neg P)\).}

Il connettivo principale ora è \(\neg\), e quindi la strada più promettente è quella di usare \(\neg I\) per ottenere \(\neg(P\land\neg P)\) a partire da \(\bot\) scaricando l'assunzione di \(P\land\neg P\).

\begin{prooftree}
  \AxiomC{\(\bot\)}
  \RuleLabel{\(\neg I\)}[\({[P\land\neg P]}_1\)]
  \UnaryInfC{\(\neg(P\land\neg P)\)}
\end{prooftree}

Dobbiamo quindi dimostrare \(\bot\) avendo a disposizione \(P\land\neg P\) come assunzione scaricabile. Essendo \(\neg E\) l'unica regola che ci permette di dedurre \(\bot\), è ragionevole pensare di usarla. In particolare, per poterla applicare è necessario avere una formula \(P\) e la sua negazione.

\begin{prooftree}
  \AxiomC{\(P\)}
  \AxiomC{\(\neg P\)}
  \RuleLabel{\(\neg E\)}
  \BinaryInfC{\(\bot\)}
  \RuleLabel{\(\neg I\)}[\({[P\land\neg P]}_1\)]
  \UnaryInfC{\(\neg(P\land\neg P)\)}
\end{prooftree}

Possiamo notare però che entrambe queste assunzioni sono deducibili dalle regole di eliminazione per \(\land\) a partire da \(P\land\neg P\), che sappiamo essere un'assunzione scaricabile, e quindi possiamo usarla come premessa per dedurre sia \(P\) che \(\neg P\).

\begin{prooftree}
  \AxiomC{\({[P\land\neg P]}_1\)}
  \RuleLabel{\(\land E^l\)}
  \UnaryInfC{\(P\)}
  \AxiomC{\({[P\land\neg P]}_1\)}
  \RuleLabel{\(\land E^r\)}
  \UnaryInfC{\(\neg P\)}
  \RuleLabel{\(\neg E\)}
  \BinaryInfC{\(\bot\)}
  \RuleLabel{\(\neg I\)}[\({[P\land\neg P]}_1\)]
  \UnaryInfC{\(\neg(P\land\neg P)\)}
\end{prooftree}

Questo esempio ci mostra come un assunzione scaricabile possa essere usata più volte come premessa per dedurre formule diverse, a patto che la regola che la scarica sia radice di un sottoalbero che contiene tutte le deduzioni che la usano come premessa. In questo caso, essendo che \(\neg I\) è la radice di tutto l'albero, possiamo usare \(P\land\neg P\) liberamente, ma se tale regola fosse stata un altro nodo interno, avremmo potuto scaricarla solo nei rami che discendevano da quel nodo, e non in altri rami.

\begin{note}{}
  Le deduzioni che abbiamo visto fino a ora sono tutte deduzioni di tautologie note. Un altro modo per farci guidare nella scelta delle regole da applicare è quello di pensare alla validità semantica della formula che vogliamo dimostrare o a eventuali equivalenze, per capire da quale assunzione \q{estrarre} l'informazione che ci manca per la dimostrazione.

  Formalmente questo processo non è corretto, poiché assume che il nostro calcolo inferenziale catturi la conseguenza logica, ovvero abbia le proprietà di correttezza e completezza, le quali, senza una dimostrazione formale, non possiamo dare per scontate.

  Per quanto riguarda il calcolo di deduzione naturale però, come dimostreremo in seguito, è possibile dimostrare che esso è corretto e completo, e quindi possiamo usare questo tipo di ragionamento per guidarci nella scelta delle regole da applicare, avendo la certezza che se una formula è valida, allora sarà dimostrabile, e viceversa.
\end{note}

\subsubsection*{Dimostrazione di \(P \lor \neg P\).}\label{subsubsec:excluded_middle_dem}

In questo caso, il connettivo principale è \(\lor\), e quindi un idea ragionevole è quella di iniziare \(\lor I\) per ottenere \(P \lor \neg P\) a partire da \(P\) o \(\neg P\).
Per ora iniziamo a ottenerlo assumendo \(P\) e usando \(\lor I^l\).

\begin{prooftree}
  \AxiomC{\(P\)}
  \RuleLabel{\(\lor I^l\)}
  \UnaryInfC{\(P \lor \neg P\)}
\end{prooftree}

Non avendo alcuna assunzione scaricabile, la strada più promettente è quella di usare \(\RAA\) per ottenere \(P\) a partire da \(\bot\) scaricando l'assunzione di \(\neg P\).

\begin{prooftree}
  \AxiomC{\(\bot\)}
  \RuleLabel{\(\RAA_1\)}[\({[\neg P]}_1\)]
  \UnaryInfC{\(P\)}
  \RuleLabel{\(\lor I^l\)}
  \UnaryInfC{\(P \lor \neg P\)}
\end{prooftree}

Ora però dobbiamo dimostrare \(\bot\) avendo a disposizione \(\neg P\) come assunzione scaricabile. Come già visto, l'unica regola che ci permette di dedurre \(\bot\) è \(\neg E\), ma per poterla applicare in questo caso è necessario assumere \(P\) come premessa andando a creare un circolo vizioso.

\begin{prooftree}
  \AxiomC{\(P\)}
  \AxiomC{\({[\neg P]}_1\)}
  \RuleLabel{\(\neg E\)}
  \BinaryInfC{\(\bot\)}
  \RuleLabel{\(\RAA_1\)}
  \UnaryInfC{\(P\)}
  \RuleLabel{\(\lor I^l\)}
  \UnaryInfC{\(P \lor \neg P\)}
\end{prooftree}

Dobbiamo quindi pensare a un altro modo per dedurre \(\bot\) a partire da \(\neg P\). Come detto prima, \(\lor I\) può essere usato per dedurre \(P \lor \neg P\) a partire sia da \(P\) che da \(\neg P\). Possiamo quindi dedurre \(P \lor \neg P\) a partire da \(\neg P\) usando \(\lor I^r\), in modo da poter dedurre \(\bot\) assumendo \(\neg(P\lor\neg P)\).

\begin{prooftree}
  \AxiomC{\({[\neg P]}_1\)}
  \RuleLabel{\(\lor I^r\)}
  \UnaryInfC{\(P \lor \neg P\)}
  \AxiomC{\(\neg(P\lor\neg P)\)}
  \RuleLabel{\(\neg E\)}
  \BinaryInfC{\(\bot\)}
  \RuleLabel{\(\RAA_1\)}
  \UnaryInfC{\(P\)}
  \RuleLabel{\(\lor I^l\)}
  \UnaryInfC{\(P \lor \neg P\)}
\end{prooftree}

A questo punto abbiamo dimostrato \(\neg (P\lor \neg P)\vdash P \lor \neg P\). Per poter arrivare al teorema, ovvero \(\vdash P \lor \neg P\), dobbiamo scaricare l'assunzione di \(\neg (P\lor \neg P)\). Per farlo possiamo pensare di riutilizzarla, in modo da riottenere \(\bot\) usando \(\neg E\), e a quel punto scaricarla con \(\RAA\) per ottenere \(P \lor \neg P\).

\begin{prooftree}
  \AxiomC{\({[\neg P]}_1\)}
  \RuleLabel{\(\lor I^r\)}
  \UnaryInfC{\(P \lor \neg P\)}
  \AxiomC{\({[\neg(P\lor\neg P)]}_2\)}
  \RuleLabel{\(\neg E\)}
  \BinaryInfC{\(\bot\)}
  \RuleLabel{\(\RAA_1\)}
  \UnaryInfC{\(P\)}
  \RuleLabel{\(\lor I^l\)}
  \UnaryInfC{\(P \lor \neg P\)}
  \AxiomC{\({[\neg(P\lor\neg P)]}_2\)}
  \RuleLabel{\(\neg E\)}
  \BinaryInfC{\(\bot\)}
  \RuleLabel{\(\RAA_2\)}
  \UnaryInfC{\(P \lor \neg P\)}
\end{prooftree}

Questo ci mostra che nonostante la regola di deduzione finale sia \(\RAA\) anche partendo da un altra regola di deduzione (\(\lor I\)) siamo riusciti a dimostrare lo stesso \(P \lor \neg P\).

\subsubsection*{Dimostrazione di \((P \to Q) \leftrightarrow (\neg P \lor Q)\).}\label{subsubsec:material_implication_dem}

Come abbiamo già visto nell'\Namedcref{ex:notable_tautologies}, \((P \to Q) \leftrightarrow (\neg P \lor Q)\) è una tautologia, quindi ci aspettiamo di poterla dimostrare. Come abbiamo già visto per dimostrare un equivalenza basta dimostrare entrambe le implicazioni che la compongono per poi concludere usando \(\land I\).

Andiamo quindi a dimostrare separatamente \((P \to Q) \to (\neg P \lor Q)\) e \((\neg P \lor Q) \to (P \to Q)\). Banalmente però, sfruttando \(\to I\) possiamo passare al dover dimostrare i sequenti \(P \to Q \vdash \neg P \lor Q\) e \(\neg P \lor Q \vdash P \to Q\), il che ci permette di non doverci preoccupare di scaricare l'assunzione di \(P \to Q\) e \(\neg P \lor Q\) rispettivamente, e quindi di poterle usare liberamente come assunzioni.

\paragraph{Dimostrazione di \(P \to Q \vdash \neg P \lor Q\).}

Dalla dimostrazione precedente, abbiamo visto che un modo per dimostrare formule con \(\lor\) come connettivo principale è quello di usare \(\RAA\).

\begin{prooftree}
  \AxiomC{\(\bot\)}
  \RuleLabel{\(\RAA_1\)}[\({[\neg(\neg P \lor Q)]}_1\)]
  \UnaryInfC{\(\neg P \lor Q\)}
\end{prooftree}

Come già detto numerose volte, per dimostrare \(\bot\) l'idea più promettente è quella di usare \(\neg E\), per la quale ci servono due formule, una la negazione dell'altra. Dall'assunzione di \(P \to Q\) possiamo dedurre \(Q\) assumendo \(P\). Quindi speriamo di poter usare \(\neg E\) con \(Q\) e \(\neg Q\).

\begin{prooftree}
  \AxiomC{\(P\)}
  \AxiomC{\(P \to Q\)}
  \RuleLabel{\(\to E\)}
  \BinaryInfC{\(Q\)}
  \AxiomC{\(\neg Q\)}
  \RuleLabel{\(\neg E\)}
  \BinaryInfC{\(\bot\)}
  \RuleLabel{\(\RAA_1\)}[\({[\neg(\neg P \lor Q)]}_1\)]
  \UnaryInfC{\(\neg P \lor Q\)}
\end{prooftree}

A questo punto dobbiamo scaricare o dimostrare sia \(P\) che \(\neg Q\), poiché \(P \to Q\) non è necessario che venga scaricata.
Dall'assunzione scaricabile di \(\neg(\neg P \lor Q)\) però, dall'equivalenza semantica con \(P \land \neg Q\), ci aspettiamo di poter dedurre sia \(P\) che \(\neg Q\) da \(\land E\).

Iniziamo con \(\neg Q\). Per poter introdurre una negazione, possiamo utilizzare \(\neg I\), che ci permette di dedurre \(\neg Q\) a partire da \(\bot\) scaricando l'assunzione di \(Q\).

\begin{prooftree}
  \AxiomC{\(P\)}
  \AxiomC{\(P \to Q\)}
  \RuleLabel{\(\to E\)}
  \BinaryInfC{\(Q\)}
  \AxiomC{\(\bot\)}
  \RuleLabel{\(\neg I_2\)}[\({[Q]}_2\)]
  \UnaryInfC{\(\neg Q\)}
  \RuleLabel{\(\neg E\)}
  \BinaryInfC{\(\bot\)}
  \RuleLabel{\(\RAA_1\)}[\({[\neg(\neg P \lor Q)]}_1\)]
  \UnaryInfC{\(\neg P \lor Q\)}
\end{prooftree}

A questo punto per dimostrare questo \(\bot\) dobbiamo introdurre una contraddizione. Avendo a disposizione \(Q\) e \(\neg(\neg P \lor Q)\) possiamo pensare di questa seconda insieme a \(\neg P \lor Q\) per dedurre una contraddizione usando \(\neg E\).

\begin{prooftree}
  \AxiomC{\(P\)}
  \AxiomC{\(P \to Q\)}
  \RuleLabel{\(\to E\)}
  \BinaryInfC{\(Q\)}
  \AxiomC{\({[\neg(\neg P \lor Q)]}_1\)}
  \AxiomC{\(\neg P \lor Q\)}
  \RuleLabel{\(\neg E\)}
  \BinaryInfC{\(\bot\)}
  \RuleLabel{\(\neg I_2\)}[\({[Q]}_2\)]
  \UnaryInfC{\(\neg Q\)}
  \RuleLabel{\(\neg E\)}
  \BinaryInfC{\(\bot\)}
  \RuleLabel{\(\RAA_1\)}
  \UnaryInfC{\(\neg P \lor Q\)}
\end{prooftree}

Ma \(\neg P \lor Q\) è deducibile da \(Q\) tramite \(\lor I^r\), permettendoci di chiudere questo ramo di deduzione.

\begin{prooftree}
  \AxiomC{\(P\)}
  \AxiomC{\(P \to Q\)}
  \RuleLabel{\(\to E\)}
  \BinaryInfC{\(Q\)}
  \AxiomC{\({[\neg(\neg P \lor Q)]}_1\)}
  \AxiomC{\({[Q]}_2\)}
  \UnaryInfC{\(\neg P \lor Q\)}
  \RuleLabel{\(\neg E\)}
  \BinaryInfC{\(\bot\)}
  \RuleLabel{\(\neg I_2\)}
  \UnaryInfC{\(\neg Q\)}
  \RuleLabel{\(\neg E\)}
  \BinaryInfC{\(\bot\)}
  \RuleLabel{\(\RAA_1\)}
  \UnaryInfC{\(\neg P \lor Q\)}
\end{prooftree}

Ci rimane quindi solo da dimostrare \(P\). Per farlo però possiamo applicare un ragionamento simile a quello fatto per \(\neg Q\), ma utilizzando \(\RAA\) invece di \(\neg I\) per dedurre \(P\) a partire da \(\bot\) scaricando l'assunzione di \(\neg P\).

\begin{prooftree}
  \AxiomC{\({[\neg(\neg P \lor Q)]}_1\)}
  \AxiomC{\({[\neg P]}_3\)}
  \RuleLabel{\(\lor I^r\)}
  \UnaryInfC{\(\neg P \lor Q\)}
  \RuleLabel{\(\neg E\)}
  \BinaryInfC{\(\bot\)}
  \RuleLabel{\(\RAA_3\)}
  \UnaryInfC{\(P\)}
  \AxiomC{\(P \to Q\)}
  \RuleLabel{\(\to E\)}
  \BinaryInfC{\(Q\)}
  \AxiomC{\({[\neg(\neg P \lor Q)]}_1\)}
  \AxiomC{\({[Q]}_2\)}
  \UnaryInfC{\(\neg P \lor Q\)}
  \RuleLabel{\(\neg E\)}
  \BinaryInfC{\(\bot\)}
  \RuleLabel{\(\neg I_2\)}
  \UnaryInfC{\(\neg Q\)}
  \RuleLabel{\(\neg E\)}
  \BinaryInfC{\(\bot\)}
  \RuleLabel{\(\RAA_1\)}
  \UnaryInfC{\(\neg P \lor Q\)}
\end{prooftree}

L'unica assunzione non scaricata che ci rimane è \(P \to Q\), e quindi abbiamo dimostrato \(P \to Q \vdash \neg P \lor Q\).

\paragraph{Dimostrazione di \(\neg P \lor Q \vdash P \to Q\).}

Ancora una volta, per dimostrare \(P \to Q\) è ragionevole pensare di usare \(\to I\) per ottenere \(P \to Q\) a partire da \(Q\) scaricando l'assunzione di \(P\).

\begin{prooftree}
  \AxiomC{\(Q\)}
  \RuleLabel{\(\to I_1\)}[\({[P]}_1\)]
  \UnaryInfC{\(P \to Q\)}
\end{prooftree}

Ora però dobbiamo dimostrare \(Q\) avendo a disposizione \(\neg P \lor Q\) e \(P\) come assunzioni. Essendoci un \(\lor\) che contiene \(Q\) come operando, è possibile pensare di utilizzare \(\lor E\) per dedurre \(Q\) a partire da \(\neg P \lor Q\), ma per poterlo applicare è necessario dimostrare \(Q\) sia assumendo \(\neg P\) che assumendo \(Q\).

\begin{prooftree}
  \AxiomC{\(\neg P \lor Q\)}
  \AxiomC{\(Q\)}
  \AxiomC{\(Q\)}
  \RuleLabel{\(\lor E_2\)}[\({[\neg P]}_2,\ {[Q]}_2\)]
  \TrinaryInfC{\(Q\)}
  \RuleLabel{\(\to I_1\)}[\({[P]}_1\)]
  \UnaryInfC{\(P \to Q\)}
\end{prooftree}

Per quanto riguarda dimostrare \(Q\) assumendo \(Q\), questo è immediato.
\begin{prooftree}
  \AxiomC{\(\neg P \lor Q\)}
  \AxiomC{\({[Q]}_2\)}
  \UnaryInfC{\(Q\)}
  \AxiomC{\(Q\)}
  \RuleLabel{\(\lor E_2\)}[\({[\neg P]}_2\)]
  \TrinaryInfC{\(Q\)}
  \RuleLabel{\(\to I_1\)}[\({[P]}_1\)]
  \UnaryInfC{\(P \to Q\)}
\end{prooftree}

Ora rimane da dimostrare \(Q\) assumendo \(\neg P\). Per farlo, possiamo utilizzare \(\bot E\), poiché tra le assunzioni scaricabili abbiamo sia \(P\) che \(\neg P\), e quindi possiamo dedurre \(\bot\) usando \(\neg E\), per poi dedurre \(Q\) da \(\bot\).
\begin{prooftree}
  \AxiomC{\(\neg P \lor Q\)}
  \AxiomC{\({[Q]}_2\)}
  \UnaryInfC{\(Q\)}
  \AxiomC{\({[\neg P]}_2\)}
  \AxiomC{\({[P]}_1\)}
  \RuleLabel{\(\neg E\)}
  \BinaryInfC{\(\bot\)}
  \RuleLabel{\(\bot E\)}
  \UnaryInfC{\(Q\)}
  \RuleLabel{\(\lor E_2\)}
  \TrinaryInfC{\(Q\)}
  \RuleLabel{\(\to I_1\)}
  \UnaryInfC{\(P \to Q\)}
\end{prooftree}

In questo esempio, abbiamo potuto osservare come anche le regole di inferenza più controintuitive e complesse, come \(\bot E\) e \(\lor E\) possano trovare applicazione.

\subsection{Proprietà del calcolo deduttivo}

Abbiamo visto come il calcolo di deduzione naturale sia un sistema formale che ci permette sintatticamente di ottenere nuove formule a partire da un insieme di premesse, ma per poterlo usare in modo efficace è necessario assicurarsi che esso sia in grado di catturare la conseguenza logica.

In questo senso, andremo a dimostrare varie proprietà del calcolo di deduzione naturale, che lo mettono in relazione con la conseguenza logica, per fare in modo da garantire che questa tecnica di manipolazione sintattica riesca a mantenere valore semantico, e quindi che se una formula è dimostrabile, allora è anche valida, e viceversa.

Iniziamo con una proprietà strettamente relativa al calcolo che ci permette di ampliare arbitrariamente l'insieme di premesse da cui dedurre una formula, senza perdere la possibilità di dedurla.

\begin{theorem}[namedlabel={thm:monotonicity}{Monotonicità}]{Monotonicità di \(\vdash\)}
  Dato un insieme di formule \(\Gamma\) e due formule \(\phi\) e \(\psi\), si ha che
  \[\Gamma \vdash \phi \implies \Gamma \cup \{\psi\} \vdash \phi\]
\end{theorem}

È importante notare che questa proprietà risulta solo rafforzata nel caso vengano aggiunte formule che rendono \(\Gamma\) inconsistente, poiché in questo caso sarà possibile dedurre qualsiasi formula a partire da \(\Gamma\) mediante \(\bot E\). Questa è una caratteristica tipica della logica classica, e non è condivisa da tutte le logiche, come ad esempio le logiche paraconsistenti, in cui l'inconsistenza non porta alla trivialità, ma viene in qualche modo contenuta permettendo di parlare del resto dell'insieme di formule anche in presenza di contraddizioni.

\begin{note}{}
  Questa proprietà di generalizza in \(\Gamma \vdash \phi \implies \Gamma \cup \Gamma' \vdash \phi\), applicandola più volte. La dimostrazione di questa versione più generale però segue direttamente da quella della versione più semplice, e quindi non è necessario farla esplicitamente.
\end{note}

\begin{proof}
  Per definizione di deduzione, \(\Gamma \vdash \phi\) se e solo se esiste un albero di derivazione ben formato radice etichettata con \(\phi\) e le cui foglie o sono etichettate con formule di \(\Gamma\) o sono assunzioni scaricabili. 

  In tale definizione non c'è alcun obbligo di usare ogni formula di \(\Gamma\) come etichetta di una foglia. Pertanto, aggiungendo all'insieme delle premesse \(\Gamma\) una nuova formula \(\psi\) non usata nell'albero di derivazione, lo stesso identico albero di derivazione rimane valido per dimostrare \(\Gamma \cup \{\psi\} \vdash \phi\).
\end{proof}

Andiamo quindi a vedere una proprietà analoga a quella vista per la conseguenza logica nel \Namedcref{thm:semantic_deduction}, il vero e proprio \Namedonlyhref{thm:deduction}.

\begin{theorem}[namedlabel={thm:deduction}{Teorema di Deduzione}]{Teorema di Deduzione}
  Dato un insieme di formule \(\Gamma\) e due formule \(\phi\) e \(\psi\), si ha che:
  \[\Gamma, \psi \vdash \phi \iff \Gamma \vdash \psi \to \phi\]
\end{theorem}

\begin{proof}
  Andiamo a dimostrare entrambi i versi.
  
  Iniziamo con \(\Gamma, \psi \vdash \phi \implies \Gamma \vdash \psi \to \phi\).
  Per definizione di deduzione, esiste un albero che dimostra \(\Gamma, \psi \vdash \phi\). Allora è sufficiente applicare la regola \(\to I\) in radice di tale albero, scaricando l'assunzione \(\psi\) e ottenendo così un albero che dimostra \(\Gamma \vdash \psi \to \phi\).

  Nell'altro verso, \(\Gamma \vdash \psi \to \phi \implies \Gamma, \psi \vdash \phi\), l'idea è duale.
  Per definizione di deduzione, esiste un albero che dimostra \(\Gamma \vdash \psi \to \phi\). A questo punto, aggiungendo l'assunzione \(\psi\) all'albero, possiamo applicare la regola \(\to E\) per ottenere un albero che dimostra \(\Gamma, \psi \vdash \phi\).
\end{proof}

Già da questi primi due teoremi è possibile notare un forte parallelismo tra la conseguenza logica e la deduzione, la quale verrà ulteriormente rafforzata dai prossimi teoremi, per poi essere dimostrata in modo definitivo con i teoremi di correttezza e completezza.

\begin{theorem}{Regola del taglio}
  Dato un insieme di formule \(\Gamma\) e due formule \(\phi\) e \(\psi\), si ha che:
  
  Se \(\Gamma, \psi \vdash \phi\) e \(\Gamma' \vdash \psi\) allora \(\Gamma \cup \Gamma' \vdash \phi\)
\end{theorem}

Intuitivamente, questa regola ci dice se abbiamo un albero di derivazione che dimostra un'assunzione, possiamo sostituire a ogni foglia etichettata con quell'assunzione l'albero di derivazione che dimostra quell'assunzione, mantenendo una derivazione valida.


\begin{proof}
  Possiamo formalmente dimostrare questa regola tramite il teorema di deduzione e la monotònicità di \(\vdash\).
  Infatti per monotònicità valgono sia \(\Gamma, \Gamma', \psi \vdash \phi\), che \(\Gamma \cup \Gamma'\vdash \psi\).
  A questo punto, usando il teorema di deduzione, otteniamo \(\Gamma\cup \Gamma' \vdash \psi \to \phi\).

  A questo punto, possiamo costruire un albero di derivazione con \(\phi\) in radice dedotta tramite \(\to E\), i cui due sottoalberi sono l'albero relativo a \(\Gamma \cup \Gamma' \vdash \psi\) e l'albero relativo a \(\Gamma\cup \Gamma' \vdash \psi \to \phi\).
\end{proof}

Prima di spostarci a dimostrare le proprietà di correttezza e completezza, è utile introdurre formalmente l'equivalente sintattico del concetto di insieme insoddisfacibile, ovvero quello di insieme inconsistente.

\begin{definition}{Inconsistenza}
  Un insieme di formule \(\Gamma\) è \keyword{inconsistente} se \(\Gamma \vdash \bot\).

  Equivalentemente, \(\Gamma\) è inconsistente, per definizione, se \(\Gamma \vdash \phi\), \(\forall \phi \in \PL\).
\end{definition}

Vedere che le due definizioni sono equivalenti è immediato.
Infatti, se \(\Gamma \vdash \bot\), possiamo applicare \(\bot E\) alla derivazione di \(\bot\), ottenendo così \(\Gamma \vdash \phi\) per ogni formula \(\phi\). Viceversa, se \(\Gamma \vdash \phi\) per ogni formula \(\phi\), allora in particolare si ha che \(\Gamma \vdash \bot\), poiché \(\bot\) è una formula di queste \(\phi\).

Veniamo dunque alle proprietà di correttezza e completezza, che intuitivamente esprimono relazioni di inclusione tra la conseguenza logica e la deduzione. 

\begin{definition}{Completezza}
  Un sistema di deduzione è \keyword{completo} se, per ogni insieme di formule \(\Gamma\) e ogni formula \(\phi\), si ha che:
  \[\Gamma \models \phi \implies \Gamma \vdash \phi\]
\end{definition}

\begin{definition}{Correttezza}
  Un sistema di deduzione è \keyword{corretto} se, per ogni insieme di formule \(\Gamma\) e ogni formula \(\phi\), si ha che:
  \[\Gamma \vdash \phi \implies \Gamma \models \phi\]
\end{definition}

Essendo che sia \(\vdash \subseteq 2^{\PL}\times\PL\) che \(\models \subseteq 2^{\PL}\times\PL\), è possibile pensare a queste due proprietà come a delle relazioni di inclusione tra i due insiemi. In particolare, diremo che il calcolo è corretto se \(\vdash \subseteq \models\), e completo se \(\models \subseteq \vdash\).

Intuitivamente, la completezza di un calcolo ci dice che tutto ciò che è conseguenza logica è anche deducibile, e quindi che il calcolo è in grado di catturare tutta la conseguenza logica. La correttezza invece ci dice che tutto ciò che è deducibile è anche conseguenza logica, e quindi che il calcolo non permette di dedurre formule che non sono conseguenza logica.

Queste due proprietà prese in isolamento però non ci dicono molto, poiché è possibile soddisfarle in modo triviale. Per la completezza, ad esempio, è sufficiente prendere come relazione di deduzione l'intero insieme \(2^{\PL}\times\PL\), in modo da poter dedurre qualsiasi formula a partire da qualsiasi insieme di formule, e quindi soddisfare la condizione di completezza, ma non quella di correttezza. Per la correttezza invece è sufficiente prendere come relazione di deduzione l'insieme vuoto, in modo da non poter dedurre nessuna formula a partire da nessun insieme di formule, e quindi soddisfare la condizione di correttezza, ma non quella di completezza.

Prima di poter dimostrare però che il calcolo di deduzione naturale è corretto e completo, sarà necessario introdurre alcune definizioni e risultati intermedi, che ci permetteranno di dimostrare più facilmente queste due proprietà.

Iniziamo col definire un concetto che ci aiuterà nella dimostrazione di correttezza.

\begin{definition}{Denotazione}
  Data una formula \(\phi\), la \keyword{denotazione} di \(\phi\) è l'insieme di tutti i modelli che soddisfano \(\phi\), ovvero:
  \[\den{\phi} = \set{m \in 2^{\AP}}[ m \models \phi]\]

  Dato un insieme di formule \(\Gamma\), la \keyword{denotazione} di \(\Gamma\) è l'insieme di tutti i modelli che soddisfano tutte le formule di \(\Gamma\), ovvero:
  \[\den{\Gamma} = \set{m \in 2^{\AP}}[ m \models \Gamma] = \bigcap_{\phi \in \Gamma} \den{\phi}\]
\end{definition}

Possiamo dunque procedere alla dimostrazione della correttezza del calcolo di deduzione naturale, che intuitivamente ci dice che tutto ciò che è provabile è semanticamente fondato, e che quindi il calcolo di deduzione naturale ci permette di dedurre solo formule che sono conseguenza logica delle premesse.

\begin{theorem}{Correttezza del calcolo di deduzione naturale}
  Per ogni insieme di formule \(\Gamma\) e ogni formula \(\phi\) si ha che:
  \[\Gamma \vdash \phi \implies \Gamma \models \phi\]
\end{theorem}

\begin{proof}

  Per definizione di \(\Gamma\vdash\phi\), esiste un albero di deduzione con radice \(\phi\) e con foglie che sono o assunzioni scaricate o formule di \(\Gamma\).

  Dunque dimostrare la tesi equivale a dimostrare che, per ogni albero di deduzione in deduzione naturale per \(\phi\) da \(\Gamma\), \(\Gamma\models\phi\).

  Possiamo facilmente vederlo per induzione sulla struttura (o altezza) dell'albero di deduzione.
  
  Il caso base è una deduzione di altezza \(0\), ovvero \(\phi \in \Gamma\). In questo caso, \(\Gamma\models\phi\) è ovviamente verificato.

  A questo punto sia l'altezza maggiore di \(0\). In questo caso avremo che l'albero avrà in radice \(\phi\) dedotta da una certa regola \(\Rule \in \Rules\) e con figli \(\psi_1,\ldots,\psi_k \in \PL\), ognuno di questi a sua volta radice di un sottoalbero di altezza strettamente minore di quella dell'albero di deduzione di \(\phi\). Per ipotesi induttiva, \(\Gamma\models\psi_i\) per ogni \(i=1,\ldots,k\).

  A questo punto quindi ci basta mostrare con un analisi per casi che, per ogni regola di deduzione naturale, se \(\set{\psi_1,\ldots,\psi_k} \vdash_{\Rule} \phi\) allora \(\set{\psi_1,\ldots,\psi_k} \models_{\Rule} \phi\).

  \begin{description}
    \item[\(\to E\).] Abbiamo da ipotesi induttiva che \(\Gamma\models\psi\) e \(\Gamma\models\psi\to\phi\).
      Ovvero, ogni modello che soddisfa \(\Gamma\) soddisfa anche \(\psi\) e \(\psi\to\phi\).
      
      Dalla semantica dell'implicazione, abbiamo che, se un modello \(m\) soddisfa \(\psi\to\phi\) allora o non soddisfa \(\psi\) o, se la soddisfa, allora soddisfa \(\phi\). Ma essendo che \(m\) soddisfa \(\psi\), allora \(m\) deve soddisfare anche \(\phi\). Dunque ogni modello che soddisfa \(\Gamma\) soddisfa anche \(\phi\), e quindi \(\Gamma\models\phi\).
    \item[\(\neg E\).] Abbiamo da ipotesi induttiva che \(\Gamma\models\neg\psi\) e \(\Gamma\models\psi\). Dunque ogni modello che soddisfa \(\Gamma\) soddisfa anche \(\neg\psi\) e \(\psi\). Per semantica del \(\neg\), allora \(m\) non soddisfa \(\psi\), ma abbiamo detto che \(m\) soddisfa \(\psi\), assurdo. Dunque non esiste alcun modello che soddisfa \(\Gamma\), e dunque \(\Gamma\) è insoddisfacibile.

    Una definizione alternativa di conseguenza logica che avevamo dato era che \(\Gamma\models\phi\) se \(\Gamma\cup\set{\neg\phi}\) è insoddisfacibile. Ma allora \(\Gamma\models\bot\) se \(\Gamma\cup\set{\neg \bot}\) è insoddisfacibile. Ma \(\neg\bot\) è tautologicamente vero, dunque \(\Gamma\cup\set{\neg \bot}\) è insoddisfacibile se e solo se \(\Gamma\) è insoddisfacibile. Cosa che abbiamo appena dimostrato nel caso in cui \(\Gamma\models\neg\psi\) e \(\Gamma\models\psi\). Dunque \(\Gamma\models\bot\).

    \item[\(\RAA\).] Abbiamo da ipotesi induttiva che \(\Gamma,\neg\phi\models\bot\). Ovvero \(\Gamma\cup\set{\neg\phi}\) è insoddisfacibile. Dunque \(\Gamma\models\phi\).
    \item[\(\lor E\).] Abbiamo da ipotesi induttiva che \(\Gamma,\psi_1\models\phi\), \(\Gamma,\psi_2\models\phi\) e \(\Gamma\models\psi_1\lor\psi_2\).
    
    Possiamo uniformare le assunzioni usando il teorema di deduzione semantica, ottenendo \(\Gamma\models\psi_1\to\phi\) e \(\Gamma\models\psi_2\to\phi\). Dunque ogni modello che soddisfa \(\Gamma\) soddisfa anche \(\psi_1\to\phi\), \(\psi_2\to\phi\) e \(\psi_1\lor\psi_2\). Ma essendo che \(m\) soddisfa \(\psi_1\lor\psi_2\), allora \(m\) soddisfa almeno una tra \(\psi_1\) e \(\psi_2\). Se \(m\) soddisfa \(\psi_1\), soddisfando \(\psi_1\to\phi\) soddisfa anche \(\phi\). Se \(m\) soddisfa \(\psi_2\), soddisfacendo \(\psi_2\to\phi\) soddisfa anche \(\phi\). Dunque \(m\) soddisfa \(\phi\) in ogni caso.
    
    Alternativamente, possiamo dimostrare questo punto usando \(\den{\phi}\) e \(\den{\Gamma}\).

    Essendo che \(\Gamma\models \psi_1 \lor \psi_2\), ogni modello che soddisfa \(\Gamma\) soddisfa anche \(\psi_1\lor\psi_2\). Dunque \(\forall m \in \den{\Gamma}, m \in \den{\psi_1\lor\psi_2}\). Alternativamente, \(\den{\Gamma}\subseteq \den{\psi_1\lor\psi_2}\). Con lo stesso ragionamento, dalla semantica della disgiunzione, \(\forall m \in \den{\psi_1\lor\psi_2}, m \in \den{\psi_1}\cup \den{\psi_2}\). Inoltre, abbiamo che \(\den{\Gamma}\subseteq\den{\psi_1\to\phi}\) e \(\den{\Gamma}\subseteq\den{\psi_2\to\phi}\), e dalla semantica dell'implicazione, si ha che \(\den{\Gamma} \cap \den{\psi_1}\subseteq\den{\phi}\) e \(\den{\Gamma}\cap\den{\psi_2}\subseteq\den{\phi}\), poiché ogni modello che soddisfa \(\Gamma\) e soddisfa \(\psi_1\) o \(\psi_2\) deve soddisfare anche \(\phi\). In particolare quindi
    \[(\den{\Gamma}\cap\den{\psi_1})\cup(\den{\Gamma}\cap\den{\psi_2})\subseteq \den{\phi}\]
    dalla distributività dell'intersezione rispetto all'unione, si ha che
    \[\den{\Gamma} \cap (\den{\psi_1}\cup\den{\psi_2})\subseteq\den{\phi}\]
    Avendo già mostrato però che \(\den{\Gamma} \subseteq \den{\psi_1}\cup\den{\psi_2}\) la loro intersezione è proprio \(\den{\Gamma}\),
    dunque \(\den{\Gamma}\subseteq \den{\phi}\), ovvero ogni modello che soddisfa \(\Gamma\) soddisfa anche \(\phi\), e dunque \(\Gamma\models\phi\).
    \item[] Per questioni di brevità non verranno trattate tutte le regole, ma solo quelle più interessanti. Gli altri casi sono analoghi a quelli già trattati, e possono essere dimostrati con un ragionamento simile a quello fatto per le regole già trattate.
  \end{description}
\end{proof}

Conclusa la dimostrazione di correttezza, possiamo passare al verso più interessante, ovvero quello di completezza, che ci dice che tutto ciò che è conseguenza logica è anche deducibile, e quindi che il calcolo di deduzione naturale è in grado di catturare tutta la conseguenza logica.

Questo risulta un po' più complesso da dimostrare, poiché, a differenza della correttezza che chiedeva semplicemente di verificare che le prove generate dal calcolo esibissero una certa proprietà, ovvero l'essere conseguenze logiche delle premesse, la completezza ci chiede di verificare che l'intero spazio delle conseguenze logiche sia coperto dal calcolo, e quindi non sarà più sufficiente verificare in maniera strutturale le prove.

Per dimostrare il teorema di correttezza infatti faremo uso di un paio di risultati intermedi, che ci permetteranno di dimostrare la completezza in modo più semplice.

\begin{lemma}[namedlabel={lem:lindenbaum}{Lindenbaum}]{Lindenbaum}
  Sia \(\Sigma\subseteq\PL\) un insieme consistente di formule. Allora è possibile costruire un insieme di formule \(\Sigma^*\) tale che \(\Sigma\subseteq\Sigma^*\) e \(\Sigma^*\) è massimalmente consistente, ovvero per ogni formula \(\psi\), o \(\psi\in\Sigma^*\) o \(\neg\psi\in\Sigma^*\).
\end{lemma}

\begin{proof}
  Consideriamo un'arbitraria enumerazione delle formule di \(\PL\), ovvero una sequenza \(\phi_1,\phi_2,\ldots\) che contiene tutte le formule di \(\PL\).
  
  Si costruisca dunque una sequenza di insiemi di formule \(\Sigma_0,\Sigma_1,\ldots\) tale che \(\Sigma_0 = \Sigma\) e che
  
  \[ \Sigma_{i+1} = \begin{cases}
                      \Sigma_i\cup\set{\phi_{i+1}} & \text{se } \Sigma_i\cup\set{\phi_{i+1}} \text{ è consistente} \\
                      \Sigma_i & \text{altrimenti}
                    \end{cases}
                    \]

  In questo modo è garantito che, se \(\Sigma\) è consistente, per induzione ogni \(\Sigma_i\) è consistente.

  Definiamo allora \(\Sigma^* = \bigcup_{i\geq 0}\Sigma_i\). Dobbiamo allora mostrare che \(\Sigma^*\) è massimalmente consistente.

  È importante notare che, nonostante \(\Sigma^*\) sia il limite di una successione di cui ogni elemento presenta una proprietà, non è detto che \(\Sigma^*\) presenti quella proprietà, e quindi è necessario dimostrarlo esplicitamente. Si pensi ad esempio a una qualsiasi successione numerica divergente. Ogni elemento di questa successione ha la proprietà di essere finito, ma il limite di questa successione è infinito, e quindi non presenta la proprietà di essere finito.

  Per dimostrare dunque che \(\Sigma^*\) è massimalmente consistente dovremo prima verificare che \(\Sigma^*\) è consistente, e poi verificare che \(\Sigma^*\) è massimale.

  Per assurdo, assumiamo che non sia consistente, ovvero \(\Sigma^*\vdash\bot\). Per definizione di deduzione, deve esistere una deduzione naturale di \(\bot\) da \(\Sigma^*\) che chiamiamo \(\pi\). Una deduzione però è un oggetto finito, dunque \(\pi\) contiene solo un numero finito di foglie, e dunque solo un numero finito di assunzioni non scaricate. Dunque dev'esistere un \(\Gamma\) finito tale che \(\Gamma\subset\Sigma^*\) e \(\Gamma\vdash\bot\). Ma essendo che ogni \(\Sigma_i\) contiene tutti i \(\Sigma_j\) con \(j<i\), posso prendere un \(i\) pari al massimo indice di una formula di \(\Gamma\) rispetto all'enumerazione scelta. Dunque \(\Gamma\subseteq\Sigma_i\), e dunque \(\Sigma_i\vdash\bot\), ovvero \(\Sigma_i\) non è consistente, assurdo.\footnote{Ciò che in questo caso ci permette di dimostrare che il limite mantiene la proprietà degli elementi della successione è il fatto che la deduzione è un oggetto finito, e quindi un'eventuale violazione della proprietà di consistenza da parte del limite deve essersi verificata già in un punto della successione, e quindi in un elemento della successione.}

  Passando dunque alla massimalità, dobbiamo dimostrare che, per ogni formula \(\psi\), o \(\psi\in\Sigma^*\) o \(\neg\psi\in\Sigma^*\) e non entrambe.

  Il fatto che non possano essere entrambe è immediato, poiché se \(\psi\) e \(\neg\psi\) fossero entrambe in \(\Sigma^*\), allora \(\Sigma^*\vdash\bot\) usando \(\neg E\), e quindi \(\Sigma^*\) sarebbe inconsistente, assurdo.

  Per vedere che almeno una delle due deve essere in \(\Sigma^*\), possiamo supporre senza perdita di generalità che \(\phi\) occorra prima di \(\neg\phi\) nell'enumerazione scelta e 
  per assurdo nessuna delle due appartenga a \(\Sigma^*\).

  Sia quindi \(\phi=\psi_j\) e \(\neg\phi=\psi_k\) con \(j<k\).
  Ma allora si ha che
  \[\phi\notin\Sigma^* \implies \psi_j \notin \Sigma_j \iff \Sigma_{j-1}\cup\set{\phi} \vdash \bot\]
  Essendo però che \(j<k\) abbiamo che \(\Sigma_{j-1}\subseteq\Sigma_{k-1}\) e quindi \(\Sigma_{k-1}\cup\set{\phi} \vdash \bot\).
  Si ha anche però che \(\psi_k\notin\Sigma^*\). Dunque
  \[\psi_k \notin \Sigma_k \iff \Sigma_{k-1}\cup\set{\neg\phi}\vdash \bot\]

  A questo punto abbiamo che \(\Sigma_{k-1}\cup\set{\phi}\vdash\bot\) e \(\Sigma_{k-1}\cup\set{\neg\phi}\vdash\bot\). Utilizzando quindi una volta \(\RAA\) e una volta \(\neg I\), si può dedurre \(\Sigma_{k-1}\vdash\neg\phi\) e \(\Sigma_{k-1}\vdash\phi\), e quindi \(\Sigma_{k-1}\vdash\bot\) usando \(\neg E\), assurdo.\footnote{È possibile dare una dimostrazione alternativa della massimalità \(\Sigma^*\) usando la formulazione di massimalità equivalente per cui \(\Sigma^*\) è massimalmente consistente se, \(\forall \phi \in \PL\setminus\Sigma^*, \Sigma^*\cup\set{\phi}\vdash\bot\).
  Per tale dimostrazione è sufficiente mostrare per assurdo che se c'è un \(\phi=\psi_i\) che non appartiene a \(\Sigma^*\) e che non rompe la consistenza, questo porta ad avere \(\Sigma_i\cup\set{\phi}\vdash\bot\), che per monotònicità implica \(\Sigma^*\cup\set{\phi}\vdash\bot\), assurdo.}
\end{proof}

A questo punto, possiamo a dimostrare una proprietà più generale, che ci permetterà di dimostrare la completezza del calcolo di deduzione naturale in modo più semplice, ovvero che un insieme di formule consistente è soddisfacibile.

Per fare ciò pero, dobbiamo andare a introdurre un concetto simile a quello di denotazione, ma in un certo senso duale, ovvero la \keyword{teoria} di un modello.

\begin{definition}{Teoria}
  Dato un modello \(m\), la \keyword{teoria} di \(m\) è l'insieme di tutte le formule che sono soddisfatte da \(m\), ovvero:
  \[\Th{m} = \set{\phi \in \PL }[ m \models \phi]\]
\end{definition}

In un certo senso, sia la denotazione che la teoria collegano il mondo sintattico a quello semantico tramite una sorta di proprietà di chiusura. Infatti, la denotazione di una formula è l'insieme di tutti i modelli che la soddisfano, e quindi è un insieme di modelli, mentre la teoria di un modello è l'insieme di tutte le formule che sono soddisfatte da quel modello, e quindi è un insieme di formule.

Andiamo dunque a vedere qualche proprietà delle teorie di modelli, che ci permetteranno di costruire un modello a partire da un insieme di formule consistente.

\begin{proposition}{Massima consistenza della teoria di un modello}
  Sia \(m\) un modello e \(\Th{m}\) la sua teoria. Allora \(\Th{m}\) è massimalmente consistente.
\end{proposition}

\begin{proof}
  Per definizione di teoria, \(\Th{m} = \set{\phi\in\PL}[m\models\phi]\). Dunque, se \(\Th{m}\) fosse inconsistente, allora \(\Th{m}\vdash\bot\), e quindi dalla correttezza del calcolo di deduzione naturale, \(\Th{m}\models\bot\), e quindi \(m\models\bot\), assurdo. Dunque \(\Th{m}\) è consistente.

  Per dimostrare che \(\Th{m}\) è massimale, dobbiamo dimostrare che, per ogni formula \(\psi\), o \(\psi\in\Th{m}\) o \(\neg\psi\in\Th{m}\) ma non entrambe. Questo però è immediato per semantica della negazione, poiché se \(m\) soddisfa \(\psi\) allora non può soddisfare \(\neg\psi\), e viceversa. Inoltre se \(m\) non soddisfa \(\psi\) allora soddisfa \(\neg\psi\).
\end{proof}

\begin{proposition}[label=prop:max-consistent-set-are-theory]{Insiemi massimalmente consistenti sono teorie}
  Sia \(\Sigma^* \subseteq \PL\) massimalmente consistente. Allora esiste un modello \(m\) tale che \(\Sigma^* = \Th{m}\).
\end{proposition}

\begin{proof}
  Dato \(\Sigma^*\) massimalmente consistente, dobbiamo costruire un modello \(m\) tale che \(\Sigma^* = \Th{m}\).
  Per fare ciò, basta definire \(m = \Sigma^*\cap\AP\), poiché essendo \(\Sigma^*\) massimalmente consistente, contiene tutte e sole le proposizioni atomiche che non portano a inconsistenza.
  Infatti le uniche proposizioni atomiche che non possono essere contenute in \(\Sigma^*\) sono quelle che, se aggiunte a \(\Sigma^*\), rendono \(\Sigma^*\) inconsistente, ovvero appaiono in qualche formula che è già presente in \(\Sigma^*\) e che le richiede con assegnamento falso per essere soddisfatte.

  A questo punto, dobbiamo mostrare che \(\Sigma^* = \Th{m}\), ovvero che per ogni formula \(\phi\), \(\phi \in \Sigma^* \iff m \models \phi\).

  Per fare ciò possiamo mostrare che \(\Sigma^*\) è chiuso per deduzione, ovvero che se \(\Sigma^*\vdash\phi\) allora \(\phi\in\Sigma^*\).
  Sia infatti per assurdo \(\phi\not\in\Sigma^*\) tale che \(\Sigma^*\vdash\phi\). Dalla massimalità di \(\Sigma^*\) si ha che \(\neg\phi\in\Sigma^*\). Dunque \(\Sigma^*\vdash\neg\phi\) banalmente, e quindi \(\Sigma^*\vdash\bot\) usando \(\neg E\), assurdo.

  A questo punto, possiamo procedere per induzione sulla struttura di \(\phi\) per mostrare che \(\phi \in \Sigma^*\) se e solo se \(m\models\phi\).

  Il caso base è una formula atomica \(p\). Se \(p\in\Sigma^*\), allora \(p\in m\) per costruzione, e quindi \(m\models p\) e viceversa.
  A questo punto, sia \(\phi\) una formula non atomica. In questo caso, \(\phi\) è costruita a partire da formule di lunghezza minore tramite l'applicazione di una regola di formazione. L'ipotesi induttiva è dunque che per ogni \(\psi\) di lunghezza minore di \(\phi\),\(\psi \in \Sigma^*\) se e solo se \(m\models\psi\).
  \begin{description}
    \item[\(\phi = \neg\psi_1\)] Dalla semantica della negazione, \(m\models \phi \iff m \not\models \psi\), e per ipotesi induttiva \(m\not\models \psi\iff \psi \notin\Sigma^*\). Dalla definizione di massimalità di \(\Sigma^*\) si ha che \(\psi\notin\Sigma^*\) se e solo se \(\neg\psi\in\Sigma^*\), ovvero \(\phi\in\Sigma^*\).
    \item[\(\phi = \psi_1 Op \psi_2\)] Tutte le dimostrazioni per operatori binari sono analoghe e si basano sull'uso della chiusura per deduzione e sulle regole di deduzione naturale specifiche dell'operatore. Vediamo come esempio il caso di \(Op=\land\).
        Dalla semantica della congiunzione, \(m\models \phi\) se e solo se \(m\models \psi_1\) e \(m\models \psi_2\). Per ipotesi induttiva, \(m\models \psi_1\) se e solo se \(\psi_1\in\Sigma^*\), e \(m\models \psi_2\) se e solo se \(\psi_2\in\Sigma^*\). Ci basta quindi mostrare che \(\psi_1,\psi_2 \in \Sigma^*\iff \psi_1\land\psi_2\in\Sigma^*\).
        Ma essendo che \(\psi_1,\psi_2 \vdash_{\land I} \psi_1\land\psi_2\) e \(\psi_1\land\psi_2\vdash_{\land E} \psi_1\) e \(\psi_1\land\psi_2\vdash_{\land E} \psi_2\), dalla chiusura di \(\Sigma^*\) per deduzione si ha la tesi.
  \end{description}
\end{proof}

\begin{lemma}[namedlabel={lem:model-existence}{Esistenza del modello}]{Esistenza del modello}
  Dato \(\Sigma\subseteq\PL\), si ha che:
  
  \(\Sigma\) è consistente \(\iff\Sigma\) è soddisfacibile.
\end{lemma}

\begin{proof}
  Per definizione di consistenza \(\Sigma \not\vdash \bot\), ovvero non esiste alcun albero di deduzione con radice \(\bot\) con foglie che sono o assunzioni scaricate o appartengono a \(\Sigma\).
  Essere soddisfacibile significa che esiste un modello \(m\) tale che \(m\models\Sigma\).

  Dunque dobbiamo andare a costruire un modello \(m\) che soddisfa \(\Sigma\) sapendo che le formule di \(\Sigma\) sono consistenti.

  Ma a questo punto, se \(\Sigma\) è consistente, allora per il \Namedcref{lem:lindenbaum} è possibile costruire un insieme di formule \(\Sigma^*\) tale che \(\Sigma\subseteq\Sigma^*\) e \(\Sigma^*\) è massimalmente consistente.
  Dunque, per la \Cref{prop:max-consistent-set-are-theory}, esiste un modello \(m\) tale che \(\Sigma^* = \Th{m}\). Dunque, essendo \(\Sigma\subseteq\Sigma^*\), allora \(\Sigma\subseteq\Th{m}\), e quindi \(m\models\Sigma\), ovvero \(m\) soddisfa tutte le formule di \(\Sigma\), e quindi \(\Sigma\) è soddisfacibile.

  Viceversa, se \(\Sigma\) è soddisfacibile, allora esiste un modello \(m\) tale che \(m\models\Sigma\). Per definizione di teoria, \(\Sigma\subseteq\Th{m}\).
  Essendo che \(\Th{m}\) è consistente, allora anche \(\Sigma\) è consistente, poiché se \(\Sigma\vdash\bot\) allora \(\Th{m}\vdash\bot\) per monotònicità, assurdo.
\end{proof}

A questo punto, siamo finalmente in grado di dimostrare la completezza del calcolo di deduzione naturale, che ci dice che tutto ciò che è conseguenza logica è anche deducibile, e quindi che il calcolo di deduzione naturale è in grado di catturare tutta la conseguenza logica.

\begin{theorem}{Completezza del calcolo di deduzione naturale}
  Per ogni insieme di formule \(\Gamma\) e ogni formula \(\phi\) si ha che:
  \[\Gamma \models \phi \implies \Gamma \vdash \phi\]
\end{theorem}

\begin{proof}
  Come già visto precedentemente, abbiamo che \(\Gamma\models\phi \iff \Gamma \cup \set{\neg\phi}\) è insoddisfacibile.

  Inoltre, possiamo facilmente vedere che \(\Gamma\vdash \phi \iff \Gamma \cup \set{\neg\phi}\) è inconsistente.

  Infatti, se \(\Gamma \cup \set{\neg\phi}\) è inconsistente allora \(\Gamma \cup \set{\neg\phi}\vdash \bot\), ovvero esiste un albero di deduzione per \(\bot\) da \(\Gamma \cup \set{\neg\phi}\). Dunque, se scarichiamo l'assunzione \(\neg\phi\) con \(\RAA\), otteniamo un albero di deduzione per \(\phi\) da \(\Gamma\), ovvero \(\Gamma\vdash\phi\).

  Viceversa, se \(\Gamma\vdash\phi\) allora esiste un albero di deduzione per \(\phi\) da \(\Gamma\). Se aggiungiamo a questo albero di deduzione un'assunzione \(\neg\phi\) e applichiamo \(\neg E\), allora otteniamo un albero di deduzione per \(\bot\) da \(\Gamma \cup \set{\neg\phi}\), ovvero \(\Gamma \cup \set{\neg\phi}\vdash\bot\), e quindi \(\Gamma \cup \set{\neg\phi}\) è inconsistente.

  Di conseguenza, \(\Gamma\models\phi \implies \Gamma\vdash\phi\) è equivalente a \(\Gamma \cup \set{\neg\phi}\) è insoddisfacibile \(\implies\) \(\Gamma \cup \set{\neg\phi}\) è inconsistente.

  A questo punto, possiamo dimostrare la proprietà di completezza per contrapposizione, ovvero che, preso un qualsiasi insieme \(\Gamma\), \(\Gamma\) consistente \(\implies\) \(\Gamma\)soddisfacibile.

  Ma per fare ciò è sufficiente il \Namedonlyhref{lem:model-existence}, che ci dice che un insieme di formule consistente è soddisfacibile. Di conseguenza, se \(\Gamma\cup\set{\neg\phi}\) è insoddisfacibile, allora è inconsistente, ovvero se \(\Gamma\models\phi\) allora \(\Gamma\vdash\phi\).
\end{proof}

Esistono infine un ultimo paio di risultati importanti riguardo il calcolo di deduzione naturale della logica proposizionale, entrambi che in un certo qual senso semplificano la trattazione di questo calcolo, uno dal fronte semantico, riducendo la grandezza dello spazio dei modelli da considerare, e uno dal fronte sintattico, permettendo di \qi{modularizzare} le dimostrazioni, e quindi di semplificarle.

\begin{theorem}{Compattezza}
  Sia \(\Gamma \subseteq \PL\) e \(\phi \in \PL\). Allora \(\Gamma \models \phi\) se e solo se esiste un sottoinsieme finito \(\Gamma_0 \subseteq \Gamma\) tale che \(\Gamma_0 \models \phi\).
\end{theorem}

Una formulazione alternativa della proprietà di correttezza in termini di soddisfacibilità è
\[\Unsat[\Gamma] \iff \exists \Gamma_0 \subseteq \Gamma, \Gamma_0 \text{ finito }: \Unsat[\Gamma_0]\]
per dualità
\[\Sat[\Gamma] \iff \forall \Gamma_0 \subseteq \Gamma, \Gamma_0 \text{ finito }: \Sat[\Gamma_0]\]
Tale teorema è di particolare importanza, poiché ci garantisce che nella verifica dell'insoddisfacibilità di un insieme di formule, è sufficiente considerare solo modelli finiti.

\begin{proof}
  Abbiamo che 
  \[\Gamma \models \phi \iff \Gamma \vdash \phi\]
  Ciò è equivalente all'esistenza di una deduzione \(\pi\) di \(\phi\) da \(\Gamma\).
  Sia allora \(\Gamma_0 = leaf(\pi)\cap \Gamma\). Ovviamente \(\Gamma_0\subseteq\Gamma\) e \(\abs{\Gamma_0}<\infty\).
  Inoltre, per costruzione di \(\Gamma_0\), \(\Gamma_0\vdash\phi\), e quindi \(\Gamma_0\models\phi\) per correttezza.
  Sia allora \(\Gamma' = \set{\psi \in \pi}[\psi \text{ foglia non scaricata di }\pi]\).
  Ma essendo \(\pi\) finito, \(\Gamma'\) è finito, e banalmente \(\Gamma' \subseteq \Gamma\) per definizione di deduzione.
  Si ha ovviamente che \(\Gamma' \vdash \phi\) per costruzione di \(\Gamma'\), e quindi \(\Gamma' \models \phi\) per correttezza.

  Nell'altro verso è ovvio per monotònicità.
  \[\Gamma'\models\phi\iff\Gamma'\vdash\phi\implies\Gamma\vdash\phi\iff\Gamma\models\phi\]
\end{proof}

Arriviamo infine a un teorema che generalizza il \Namedcref{thm:substitution}, che riguarda la sostituzione nella conseguenza logica invece che nella validità

\begin{theorem}{Invarianza della Conseguenza Logica rispetto alla Sostituzione}
  Sia \(\Gamma \subseteq \PL\), \(\phi \in \PL\) e \(\sigma:\AP \to \PL\) sostituzione. Allora,
  \[\Gamma\models\phi\implies\Gamma\sigma\models\phi\sigma\]
\end{theorem}

\begin{proof}
  Abbiamo già visto nel teorema di sostituzione che \(\models\phi \implies \models\phi\sigma\) per ogni \(\sigma\).

  Possiamo usare in questo caso il teorema di compattezza. Infatti da \(\Gamma\models\phi\) segue che esiste un sottoinsieme finito \(\Gamma_0 \subseteq \Gamma\) tale che \(\Gamma_0\models\phi\).
  Essendo \(\Gamma_0\) finito sappiamo che \(\Gamma_0\models\phi\iff\models(\bigwedge\Gamma_0)\to\phi\), e quindi \(\models((\bigwedge_{\phi\in\Gamma_0}\phi)\to\phi)\sigma\) per il teorema di sostituzione.

  Per la definizione di applicazione di una sostituzione \(\sigma\) si ha che \(((\bigwedge\Gamma_0)\to\phi)\sigma\) è equivalente a \((\bigwedge_{\phi\in\Gamma_0}\phi\sigma)\to\phi\sigma\).
  Quindi \(\Gamma_0\sigma\models\phi\sigma\) per il teorema di deduzione semantica e quindi \(\Gamma\sigma\models\phi\sigma\) per monotònicità.
\end{proof}

\begin{corollary}{Sostituzione per la Deducibilità}
  \(\Gamma\vdash\phi\implies\Gamma\sigma\vdash\phi\sigma\), con \(\sigma:\AP \to \PL\) sostituzione.
\end{corollary}
\begin{proof}
  Segue banalmente dal teorema precedente e da completezza e correttezza.
\end{proof}

Grazie a questo teorema, è possibile generare nuove deduzioni a partire da deduzioni già esistenti, semplicemente applicando una sostituzione uniformemente a tutte le formule che compaiono nella deduzione. In questo modo, è possibile \qi{modularizzare} le dimostrazioni, ovvero dimostrare prima un risultato per una formula generica, e poi ottenere il risultato per formule specifiche semplicemente applicando una sostituzione alla deduzione della formula generica.

\section{Calcolo di Risoluzione proposizionale}

Il calcolo di deduzione naturale è un calcolo di deduzione molto potente, che come da nome suggerisce, risulta naturale per la dimostrazione di teoremi, poiché le regole di deduzione si allineano molto bene con le tecniche di dimostrazione che siamo abituati a utilizzare.

Tuttavia, nonostante sia molto naturale per il ragionamento umano, il calcolo di deduzione naturale non è particolarmente adatto per l'automatizzazione, poiché la ricerca di una deduzione per una formula con un numero elevato di regole da poter applicare risulta difficile da gestire algoritmicamente.

Andiamo quindi a introdurre un calcolo di deduzione alternativo, che risulta più adatto per l'automatizzazione, poiché composto da un'unica regola di deduzione, il \keyword{calcolo di risoluzione}. La presenza di un unica regola di deduzione elimina la necessità di dover scegliere quale regola applicare, lasciando come unica scelta da fare quella di quali formule utilizzare per applicare la regola, e quindi semplificando notevolmente la ricerca di una deduzione.

Una prima versione fu proposta da J. A. Robinson nel 1965, e da allora è stata oggetto di numerose estensioni e miglioramenti.

Tale calcolo non si concentra sulla conseguenza logica, ma sull'inconsistenza. 
Nonostante ciò, per decidere \(\Gamma\models\phi\) è possibile farlo decidendo se \(\Gamma\cup\set{\neg\phi}\) è insoddisfacibile.
In questo senso la risoluzione non è un calcolo \qi{generativo}, come lo era la deduzione naturale, in cui dato un insieme di formule veniva generata la conseguenza logica.

Inoltre, tale calcolo non si applica direttamente a formule proposizionali, ma solo a formule in \keyword{forma clausale}.

\begin{definition}{Forma clausale}
  Una formula \(\phi\) è in forma clausale se \(\phi \in 2^{2^\mathcal{L}}\), dove \(\mathcal{L}\) è un insieme di letterali.

  In altre parole, una formula è in forma clausale se è un insieme di insiemi di letterali.
\end{definition}

Ovviamente, essendo insiemi, l'ordine degli elementi non conta, ogni insieme \qi{interno} non può contenere due volte lo stesso letterale, e l'insieme esterno non può contenere due volte lo stesso insieme di letterali.

Esiste inoltre una corrispondenza biunivoca tra formule in forma clausale e formule proposizionali in forma normale congiuntiva.
È possibile infatti interpretare ogni insieme interno come una disgiunzione di letterali, e l'insieme esterno come una congiunzione di tali disgiunzioni.

Inoltre possiamo notare che dato un insieme di formule in \(\CNF\) la congiunzione di tutte le formule è una formula in \(\CNF\) equivalente a tutte le formule dell'insieme, e quindi possiamo sempre rappresentare un insieme di formule \(\Gamma\subseteq\CNF\) con la formula \(\bigwedge_{\phi\in\Gamma}\phi\). E per la forma clausale chiaramente questo equivale a fare l'unione tra formule in forma clausale.

Ovviamente quindi per ogni \(\Sigma\subseteq\PL\) esiste un \(F_\Sigma\) in forma clausale equivalente a \(\Sigma\).
Questo si riduce a dire che \(\forall\phi\in\PL\) esiste un \(F_\phi\) in forma clausale equivalente a \(\phi\), e questo è vero poiché ogni formula è equivalente a una formula in \(\CNF\) e ogni formula in \(\CNF\) è equivalente a una formula in forma clausale.

Quando \(F = \emptyset\) è detta forma clausale banale ed è valida.

Questo poiché, intuitivamente, \(m\models F\) se e soltanto se \(\forall C \in F\) insieme di letterali, \(m\models C\) cioè \(m\models l\) per qualche \(l \in C\) per ogni \(C\in F\).
Quindi \(m\models F \iff \forall C \in F ,\exists l \in C\st m\models l\). Ovviamente se \(F =\emptyset\) qualsiasi \(m\) soddisfa \(F\) poiché non c'è nessuna clausola da poter non soddisfare.


È importante notare che la forma clausale vuota \(F=\emptyset\) non coincide con la clausola vuota \(C=\emptyset\), tipicamente indicata con \(\square\). Per la definizione di soddisfacibilità, \(m\not\models \square\) per ogni \(m\), e quindi \(\square\) è insoddisfacibile, mentre \(F=\emptyset\) è soddisfatta da ogni \(m\).

Essendo \(\square\) insoddisfacibile, per valutare l'insoddisfacibilità di un insieme di formule in forma clausale, è sufficiente verificare se è possibile dedurre \(\square\) da quell'insieme di formule.

\begin{example}{}
  Sia 
  \[\phi=(P\to(Q\land R))\land(Q\lor(Q\to P)) \land((\neg P \to Q)\lor(\neg Q \to R))\]
  che in \(\CNF\) è equivalente a
  \[\phi_{\CNF} = (\neg P\lor Q)\land(\neg P\lor R)\land(Q\lor \neg Q \lor P)\land(P\lor Q \lor Q \lor R)\]
  ottenibile applicando De Morgan, distributività, implicazione materiale e semplificazione di doppie negazioni.
  La forma clausale equivalente è quindi
  \[F_\phi = \set{\set{\neg P, Q},\set{\neg P, R},\set{Q, \neg Q, P},\set{P, Q, R}}\]
  Ogni clausola che contiene \(P\) e \(\neg P\) per un qualche \(P \in AP\) è tautologica, ed essendo \(\phi\land\top\equiv\phi\) è possibile rimuovere la clausola \(\set{Q, \neg Q, P}\) da \(F_\phi\) senza alterare la soddisfacibilità di \(F_\phi\).

  Da ora in poi quindi assumeremo che le forme clausali non contengano clausole banali.
\end{example}

Per andare a definire formalmente il calcolo di risoluzione, ovvero definire la sua unica regola di deduzione, anch'essa chiamata risoluzione, è necessario prima definire una proprietà di clausole che sarà la condizione di applicabilità della regola di risoluzione, ovvero la proprietà di essere in \keyword{contrasto} o in \keyword{clash}.

\begin{definition}{Clausole in contrasto}
  Siano \(C_1\) e \(C_2\) due clausole. Diremo che \(C_1\) e \(C_2\) sono in \keyword{contrasto}, o in \keyword{clash}, se esiste un \(l \in \mathcal{L}\) tale che \(l \in C_1\) e \(\tilde{l} \in C_2\) con
  
  \[\tilde{l} = \begin{cases}
                \neg P & \text{se } l = P\\
                P & \text{se } l = \neg P
              \end{cases}
  \]
\end{definition}

In altre parole, due clausole sono in clash se contengono un letterale su cui non sono concordi. Intuitivamente, per soddisfarle entrambe, è necessario che almeno una delle due \qi{rinunci} ad avere ragione su quel letterale, dovendo aver soddisfatto un letterale diverso da quello di clash. Non potendo però sapere a priori quale delle due clausole rinuncerà ad avere quel letterale soddisfatto, l'unica cosa che possiamo intuitivamente dedurre è che almeno uno dei letterali della prima clausola o almeno uno della seconda dovrà essere soddisfatto. Tale ragionamento è proprio ciò che caratterizza la regola di risoluzione

\begin{definition}{Regola di risoluzione}
  Siano \(C_1\) e \(C_2\) due clausole in contrasto su un letterale \(l\), con \(l \in C_1\) e \(\tilde{l} \in C_2\). Allora, è possibile dedurre la clausola \(R\) detta \keyword{risolvente} di \(C_1\) e \(C_2\) su \(l\), definita come:

  \[\infer[\text{Risoluzione}]{R = (C_1 \cup C_2)\setminus\set{l, \hat{l}}}{C_1 \quad C_2}\]

\end{definition}

Tale regola di deduzione segue esattamente il ragionamento intuitivo fatto in precedenza applicato alle forme clausali, in cui le clausole sono considerate in congiunzione, ovvero da soddisfare entrambe, e i letterali al loro interno in disgiunzione. L'insieme \((C_1 \cup C_2)\setminus\set{l, \hat{l}}\) è esattamente l'insieme contenente tutti i letterali delle due clausole tranne quelli in clash. Interpretandoli in disgiunzione, tale insieme corrisponde semanticamente esattamente al richiedere che ne sia soddisfatto almeno uno, che intuitivamente dovrà appartenere alla clausola che non avrà il letterale di clash soddisfatto.

Andiamo dunque a formalizzare la fondatezza semantica di tale regola.

\begin{proposition}{}
  La regola di risoluzione preserva la verità
\end{proposition}

\begin{proof}
  Formalmente, una regola di inferenza preserva la verità se ogni modello che soddisfa le premesse soddisfa anche la conclusione della regola.
  Chiaramente, il preservare la verità implica il preservare anche la conseguenza logica.\footnote{Questa cosa risulta banale nella logica proposizionale, me non sarà scontata nel prim'ordine.}

  Sia \(m\) un modello che soddisfa \(C_1\) e \(C_2\) che sono in clash su \(l\).
  Per definizione di \(l\) e \(\tilde{l}\), \(m\) soddisfa esattamente uno tra \(l\) e \(\tilde{l}\). Senza perdita di generalità, supponiamo che \(l\in C_1\) e \(\tilde{l}\in C_2\).

  Consideriamo il caso \(m\models l\). Essendo che \(m\models C_2\), allora \(\exists l' \in C_2\setminus{\tilde{l}}\) tale che \(m\models l'\). Dunque \(l' \in R\), ovvero \(m\models R\).
  Nel caso in cui \(m \models \tilde{l}\), essendo che \(m\models C_1\), allora \(\exists l' \in C_1\setminus{l}\) tale che \(m\models l'\). Dunque \(l' \in R\), ovvero \(m\models R\).
\end{proof}

Per quanto riguarda il verso opposto invece, l'unica proprietà che è preservata è la soddisfacibilità, ovvero se \(R\) risolvente di \(C_1\) e \(C_2\) su \(l\) è soddisfacibile, allora anche \(C_1\) e \(C_2\) sono soddisfacibili. Tale proprietà è ovviamente più debole della preservazione della verità, poiché l'unica garanzia che fornisce è che se esiste un modello in cui \(R\) è soddisfatto ne esiste uno, potenzialmente diverso, in cui \(C_1\) e \(C_2\) sono congiuntamente soddisfatte.

Non è possibile richiedere di più poiché \(R\) contiene un numero di proposizioni atomiche strettamente inferiore sia di \(C_1\) che di \(C_2\), e di conseguenza gli assegnamenti che soddisfano \(R\) non sono assegnamenti validi per soddisfare \(C_1\) e \(C_2\). Andiamo però a mostrare che dato \(\alpha: \AP[R] \to \set{0,1}\) tale che \(\Val{\alpha}[R] = 1\), esiste un assegnamento \(\alpha':\AP[C_1]\cup\AP[C_2] \to \set{0,1}\) tale che \(\Val{\alpha'}[C_1] = 1\) e \(\Val{\alpha'}[C_2] = 1\).

\begin{proposition}{}
  Sia \(R\) risolvente di \(C_1\) e \(C_2\) in clash su \(l\). Allora, se \(R\) è soddisfacibile, allora anche \(C_1\) e \(C_2\) sono soddisfacibili.
\end{proposition}

\begin{proof}
  Sia \(\alpha: \AP[R] \to \set{0,1}\) tale che \(\Val{\alpha}[R] = 1\).

  Dato che \(\Val{\alpha}[R] = 1\), per la semantica della forma clausale, \(\exists l' \in R\) tale che \(\Val{\alpha}[l'] = 1\).
  Per costruzione di \(R\), \(l'\) è un letterale che appartiene a \(C_1\) o a \(C_2\) o a entrambe.

  Consideriamo per primo il caso in cui \(l'\) appartiene a entrambe. In tal caso l'assegnamento \(\alpha': \AP[C_1\cup C_2]\to\set{0,1}\) definita come
  \[\alpha'(P) = \begin{cases}
                \alpha(P) & \text{se } P \in AP(R)\\
                0 & \text{se } P \not\in AP(R) 
              \end{cases}\]
              soddisfa sia \(C_1\) che \(C_2\), poiché \(l'\) appartiene a \(R\) e quindi tale letterale soddisfa sia \(C_1\) che \(C_2\).

  Consideriamo ora il caso in cui \(l'\) appartiene a \(C_1\) ma non a \(C_2\). In tal caso, l'assegnamento \(\alpha': \AP[C_1\cup C_2]\to\set{0,1}\) definito come
  \[\alpha'(P) = \begin{cases}
                \alpha(P) & \text{se } P \in AP(R)\\
                0 & \text{se } \tilde{l} = \neg P\\
                1 & \text{se } \tilde{l} = P
              \end{cases}
  \]
  soddisfa ancora entrambe, poiché \(C_1\) sarà soddisfatta da \(l'\), mentre \(C_2\) sarà soddisfatta da \(\tilde{l}\), poiché il suo assegnamento è costruito per essere soddisfatto qualunque sia la sua polarità.

  Analogamente, se \(l'\) appartiene a \(C_2\) ma non a \(C_1\), allora l'assegnamento \(\alpha': \AP[C_1\cup C_2]\to\set{0,1}\) definito come
  \[\alpha'(P) = \begin{cases}
                \alpha(P) & \text{se } P \in AP(R)\\
                0 & \text{se } l = \neg P\\
                1 & \text{se } l = P
              \end{cases}
  \]
  soddisfa ancora entrambe, per un ragionamento duale rispetto al caso precedente.
\end{proof}

Avendo ora un insieme di regole di inferenza, in questo caso composto da una regola sola, e un insieme di assunzioni, in questo caso la forma clausale, possiamo andare a definire il calcolo di deduzione di risoluzione attraverso la definizione di una deduzione.

\begin{definition}{Deduzione di risoluzione}
  Sia \(F\) una forma clausale. Una deduzione di risoluzione \(F\vdash_{\Rule }C\) da è un albero binario finito, la cui radice è etichettata con \(C\), tutte le foglie sono etichettate con clausole di \(F\) e ogni nodo interno è etichettato con il risolvente dei suoi figli.
\end{definition}

\begin{note}{}
  Essendo che in risoluzione \(\Rules = \set{\Rule}\), con \(\Rule\) la regola di risoluzione, per distinguere tale calcolo da quello di deduzione naturale indicheremo le deduzioni di risoluzione con \(\vdash_{\Rule}\) invece che con \(\vdash\). Quando sarà necessario esplicitare l'utilizzo del calcolo di deduzione naturale invece, per dissipare ogni ambiguità, verrò utilizzato \(\vdash_{ND}\).
\end{note}

Essendo che la regola di risoluzione preserva la verità, e quindi la conseguenza logica, è facile vedere per induzione strutturale la correttezza del calcolo, ovvero che date \(S,C\) forme clausali, \(S\vdash_{\Rule} C \implies S\models C\).

Sarà possibile dimostrare anche una forma più debole della completezza, legata alla soddisfacibilità, ovvero che \(S\not\vdash_R \square \implies \Sat[S]\), o equivalentemente \(\Unsat[S] \implies S\vdash_R \square\). Come già anticipato, non è possibile dimostrare la totale completezza, perché esistono conseguenze logiche di clausole che non sono in alcun modo generabili dal calcolo.

Consideriamo infatti la forma clausale \(F = \set{\set{A},\set{B}}\). Dall'interpretazione semantica data, è ovvio che \(R=\set{A,B}\) sia conseguenza logica di \(F\). Essendo però che non esistono letterali in clash in \(F\), non è possibile applicare la regola di risoluzione, impedendo di derivare alcunché.

È possibile però in parte aggirare questo problema utilizzando la definizione alternativa di conseguenza logica, notando che \(S\models C \iff S\cup\set{\neg C} \in \Unsat\).
Applicando tale riformulazione all'esempio di prima, potremo dimostrare che \(F\models R\) cercando di mostrare che \(F' = \set{\set{A}}, \set{B}, \set{\neg A, \neg B}\) è insoddisfacibile. In questo caso, essendo che \(C_1 = \set{A}\) e \(C_2 = \set{\neg A, \neg B}\) sono in clash su \(A\), è possibile applicare la regola di risoluzione per dedurre \(C_3 = \set{\neg B}\), da cui è possibile dedurre \(C_4 = \square\) con \(C_2\), che è in clash con \(C_3\) su \(B\). Dunque, \(F'\) è insoddisfacibile per correttezza, e quindi \(F\models R\) per la definizione di conseguenza logica.

Utilizzando questo tipo di riformulazione, e questa completezza \qi{debole}, è possibile risolvere il problema di conseguenza logica usando risoluzione, passando per il problema di soddisfacibilità.

Prima della dimostrazione formale di questa forma di completezza debole però, andiamo a vedere come automatizzare il calcolo di risoluzione tramite un algoritmo.

\begin{algorithmdesc}{Risoluzione}
    \begin{algorithmic}[1]
      \Statex{} \textbf{signature} \(\textsc{Resolution:} \quad [Clause] \to \Unsat \mid \Sat\)
      \Function{\(\textsc{Resolution}\)}{$S$} % chktex 46
          \State{} \(S_0\gets S\)
          \State{} \(E \gets \emptyset\)
          \Comment{\(E \subseteq S_0 \times S_0\)}
          \Repeat{}
            \State{} Pick \((C_1, C_2) \in (S_0 \times S_0)\setminus E \st \operatorname{clash}(C_1, C_2)\)
            \State{} \(R \gets \operatorname{resolve}(C_1, C_2)\)
            \If{\(R\) non banale}
              \State{} \(S_0 = S_0 \cup \set{R}\)
            \EndIf{}
            \State{} \(E \gets E \cup \set{(C_1, C_2)}\)
          \Until{\(R = \square \lor E = S_0 \times S_0 \cap \set{(C_1,C_2) \in S_0\times S_0}[\operatorname{clash}(C_1, C_2)] \)}
          \If{\(R = \square\)}
            \State{} \Return{\(\Unsat\)}
          \Else{}
            \State{} \Return{\(\Sat\)}
          \EndIf{}
      \EndFunction{}
    \end{algorithmic}
\end{algorithmdesc}

Per quanto riguarda la correttezza dell'algoritmo, questa risulta quasi scontata dal fatto che è l'esatta trasposizione di una deduzione di risoluzione. Riguardo alla terminazione invece, la cosa è meno scontata. Questo poiché la condizione di terminazione è e dev'essere valutata sull'insieme \(S_0\), il quale non è statico, ma cresce dinamicamente durante l'esecuzione. In linea di principio quindi, senza mostrare che esista un limite superiore al numero di clausole che è possibile generare mediante risoluzione, non è possibile garantire la terminazione di questo algoritmo.

Fortunatamente, tale limite superiore esiste, e ne vedremo dimostrazione formale dopo un esempio di applicazione di tale algoritmo.

\begin{example}{}
  Consideriamo l'\Cref{ex:logical_consequence} visto nel primo capitolo, formalizzato usando proposizioni atomiche per rappresentare le affermazioni legate al nome della persona coinvolta in ognuna di esse. In particolare sia:

  \begin{itemize}
    \item \(C\) la proposizione atomica che rappresenta l'affermazione \qi{Carlo è italiano}
    \item \(G\) la proposizione atomica che rappresenta l'affermazione \qi{Giovanni è spagnolo}
    \item \(R\) la proposizione atomica che rappresenta l'affermazione \qi{Rachele è tedesca}
    \item \(L\) la proposizione atomica che rappresenta l'affermazione \qi{Lucia è francese}
  \end{itemize}

  Possiamo formalizzare il problema creando l'insieme di formule \(\Gamma\):
  \begin{itemize}
    \item \((C\land \neg G)\to R\)
    \item \(R \to (L\lor G)\)
    \item \(\neg L \to C\)
    \item \(\neg G\)
  \end{itemize}

  Volevamo dimostrare che tale insieme avesse come conseguenza logica \(L\).
  Con risoluzione, prendiamo \(\Gamma \cup \set{\neg L}\) e verificare che sia insoddisfacibile.
  Dovremo trasformare quindi le formule in forma clausale passando per la \(\CNF\), e otteniamo:
  \begin{enumerate}
    \item \(\neg(C\land\neg G) \lor R \equiv \neg C \lor G \lor R\), ovvero \(\set{\neg C,G,R}\)
    \item \(\neg R \lor L \lor G\), ovvero \(\set{\neg R, L, G}\)
    \item \(L \lor C\), ovvero \(\set{L, C}\)
    \item \(\neg G\), ovvero \(\set{\neg G}\)
    \item \(\neg L\), ovvero \(\set{\neg L}\)
  \end{enumerate}

  Dobbiamo allora vedere se da queste \(5\) clausole è possibile derivare la clausola vuota \(\square\) tramite risoluzione.

  In primo luogo ho diverse scelte, prendiamo ad esempio le clausole \(2\) e \(4\) che fanno clash su \(G\). Il risolvente è \(\set{\neg R, L}\), chiamiamo questo risolvente \(6\)

  Da \(1\) e \(4\) che fanno clash su \(G\) ottengo \(\set{\neg C, R}\), chiamiamo questo risolvente \(7\).
  Da \(3\) e \(7\) otteniamo \(\set{L, R}\), chiamiamo questo risolvente \(8\).
  Da \(6\) e \(8\) otteniamo \(\set{L}\), chiamiamo questo risolvente \(9\).
  Da \(5\) e \(9\) otteniamo \(\square\).

  Ovviamente tale esecuzione non è unica. Ad esempio, si può iniziare da \(1\) e \(2\) in clash su \(R\) per ottenere \(\set{\neg C, G, L}\), che chiamiamo \(6\).

  Possiamo quindi risolvere \(6\) con \(3\) ottenendo \(\set{G, L}\), che chiamiamo \(7\).
  A questo punto risolvendo prima con \(4\), e poi il risolvente con \(5\) otteniamo \(\square\).
\end{example}

È fondamentale notare che l'applicazione della regola di risoluzione richiede la selezione di due clausole in clash su uno specifico letterale, e che quindi una singola applicazione della regola porti all'eliminazione di un solo letterale alla volta.

Non è possibile estendere banalmente la regola per permettere la risoluzione di più letterali contemporaneamente, poiché in tal caso è facile costruire controesempi in cui non è preservata la verità. Prendiamo infatti le clausole \(C_1 = \set{A, B}\) e \(C_2 = \set{\neg A, \neg B}\). Se fosse possibile risolvere entrambe le coppie di letterali in clash contemporaneamente, allora sarebbe possibile dedurre la clausola vuota \(\square\) da \(C_1\) e \(C_2\), che è insoddisfacibile, mentre invece \(C_1\) e \(C_2\) sono congiuntamente soddisfacibili da un assegnamento che assegna \(1\) ad \(A\) e \(0\) a \(B\), o viceversa. Andando ad applicare la regola di risoluzione in modo corretto, è possibile dedurre \(\set{A,\neg A}\) da \(C_1\) e \(C_2\) risolvendo su \(B\), ottenendo la clausola banale \(\set{A, \neg A}\), che è tautologica.

Fatte queste precisazioni, possiamo ora dimostrare la terminazione dell'algoritmo di risoluzione, mostrando che esiste un limite superiore al numero di clausole che tale algoritmo può generare, limitando quindi la grandezza dell'insieme \(S_0\) su cui viene valutata la condizione di terminazione.

\begin{theorem}{}
  \(\textsc{Resolution}(S)\) termina sempre e l'insieme di saturazione \(S_0\) avrà al più \(3^{\abs{\AP[S]}}\) clausole.
\end{theorem}

\begin{proof}
  Sappiamo che l'insieme di clausole in ingresso \(S\) è finito. Sia dunque \(\AP[S] = \set{P \in \AP[\phi]}[\phi \in S]\) l'insieme di proposizioni atomiche che compaiono in \(S\). Chiaramente \(\AP[S]\) è finito, poiché \(S\) è finito e ogni formula contiene un numero finito di proposizioni atomiche.

  Inoltre, avendo escluso le clausole banali, si ha che \(\forall C \in S, \abs{C}\leq\AP[S]\). Questo poiché ogni clausola è un insieme di letterali, che può essere presente o in forma positiva o in forma negativa, ma non entrambe, e quindi per ogni proposizione atomica in \(\AP[S]\) al più una delle due forme può essere presente in una clausola.

  Quindi, in ogni clausola, una data proposizione atomica \(P\) può essere presente in forma positiva, in forma negativa, o non essere presente affatto. Di conseguenza ogni clausola può essere vista come una funzione da \(\AP[S]\) a \(\set{-1,0,1}\), dove \(0\) rappresenta l'assenza di \(P\) nella clausola, \(1\) rappresenta la presenza di \(P\) in forma positiva, e \(-1\) rappresenta la presenza di \(P\) in forma negativa.

  Dunque, il numero di clausole che è possibile generare con le proposizioni atomiche in \(\AP[S]\) è al più il numero di funzioni da \(\AP[S]\) a \(\set{-1,0,1}\), che è \(3^{\abs{\AP[S]}}\).
\end{proof}

A questo punto, avendo dimostrato la correttezza e la terminazione dell'algoritmo di risoluzione, rimane da dimostrarne solo la completezza. Essendo che l'algoritmo per com'è definito genera ogni formula che il calcolo di risoluzione è in grado di generare, è sufficiente dimostrare la completezza del calcolo di risoluzione, ovvero che se \(S\) è insoddisfacibile allora \(S\vdash_R \square\).

\begin{theorem}{Completezza debole del calcolo di risoluzione}
  Se \(S\) è insoddisfacibile, allora \(S\vdash_R \square\).
\end{theorem}

\begin{proof}
  Per fare ciò, ragioniamo induttivamente sulla cardinalità di \(\AP[S]\).

  Il caso base è rappresentato da \(\abs{\AP[S]} = 0\), ovvero \(\AP[S] = \emptyset\). In questo caso è chiaro che l'unica formula che è possibile formare è la clausola vuota \(\square\), che è insoddisfacibile, e quindi \(S\vdash_R \square\).\footnote{È importante notare che \(\AP[S]=\emptyset\centernot\implies S = \emptyset\) poiché \(S\) è un insieme di insiemi di letterali. \(\AP[S]=\emptyset \implies  S = \emptyset \lor S = \set{\emptyset} = \set{\square}\). Avendo però che \(S\) è insoddisfacibile per ipotesi, siamo nella seconda situazione, e quindi \(S = \set{\square}\).}

  Per \(\abs{\AP[S]}>0\), si ha l'ipotesi induttiva che per ogni \(S'\) con \(\abs{\AP[S']} < \abs{\AP[S]}\), se \(S'\) è insoddisfacibile allora \(S'\vdash_R \square\).

  Per poter applicare tale ipotesi induttiva è quindi necessario in qualche modo rimuovere una proposizione atomica da \(S\), mostrare che il nuovo insieme di clausole è ancora insoddisfacibile, e da questa costruzione usando l'ipotesi induttiva mostrare che è possibile dedurre \(\square\) anche da \(S\).

  Iniziamo col fissare la proposizione atomica \(P\in\AP[S]\) che andremo a escludere. Costruiamo dunque gli insiemi \(S_P\) e \(S_{\neg P}\) come segue:
  \[S_P = \set{C\setminus P}[C\in S \text{ e } \neg P \in C]\]
  \[S_{\neg P} =  \set{C\setminus \neg P}[C\in S \text{ e }  P \in C]\]

  In altre parole, \(S_P\) è l'insieme di clausole che non contengono il letterale negativo di \(P\), a cui se presente è stato rimosso il letterale positivo di \(P\) e viceversa per \(S_{\neg P}\). In questo modo è chiaro che \(P\) non è presente in \(S_P\) e \(\neg P\) non è presente in \(S_{\neg P}\), e che inoltre \(\AP[S_P] \subset \AP[S]\) e \(\AP[S_{\neg P}]\subset \AP[S]\) poiché non sono state aggiunte nuove proposizioni atomiche, ma solo eventualmente rimosse.

  Di conseguenza \(\AP[S_P]<\AP[S]\) e \(\AP[S_{\neg P}]<\AP[S]\).
  Inoltre questi due insiemi non sono necessariamente disgiunti. In particolare infatti, ogni clausola in \(S\) che non conteneva ne \(P\) ne \(\neg P\) è presente sia in \(S_P\) che in \(S_{\neg P}\). In ogni caso, tutte le clausole dell'intersezione tra \(S_P\) e \(S_{\neg P}\) erano già presenti in \(S\), poiché se una clausola è stata ottenuta da \(S\) rimuovendo \(P\) o \(\neg P\), può essere ottenuta in tal modo solo in uno dei due insiemi, e quindi se è presente in entrambi, allora era già presente in \(S\) senza bisogno di rimuovere \(P\) o \(\neg P\).

  In aggiunta, tali insiemi sono anche certamente non vuoti. Supponiamo infatti per assurdo che \(S_P = \emptyset\). Allora, necessariamente, ogni clausola di \(S\) contiene \(\neg P\). Essendo però che per soddisfare una clausola è sufficiente soddisfare un suo letterale, allora ogni assegnamento che assegna \(0\) a \(P\) soddisfa tutte le clausole di \(S\), rendendo \(S\) soddisfacibile, in contraddizione con l'ipotesi che \(S\) sia insoddisfacibile. Vale ovviamente un ragionamento duale per \(S_{\neg P}\).

  Per poter applicare l'ipotesi induttiva a questi due insiemi ci rimane da mostrare che sono entrambi insoddisfacibili.
  Iniziamo ragionando sull'insoddisfacibilità di \(S_P\) supponendo per assurdo che \(S_P\) sia soddisfacibile. Allora, esiste un assegnamento \(\alpha: \AP[S_P] \to \set{0,1}\) tale che \(\Val{\alpha}[C] = 1\) per ogni \(C \in S_P\). Estendendo banalmente tale assegnamento a \(S\) con:
  \[\alpha^P(X) = \begin{cases}
    0 &, X=P\\
    \alpha(X) &, X\neq P
  \end{cases}\]
  possiamo facilmente vedere che \(\Val{\alpha^P}[S] = 1\) portando all'assurdo. Questo si ha sempre perché per soddisfare una clausola è sufficiente soddisfare un suo letterale, e ogni clausola di \(S\) o è presente in \(S_P\) e quindi soddisfatta da \(\alpha\), o contiene \(\neg P\) e quindi è soddisfatta da \(\alpha^P\) assegnando \(0\) a \(P\).
  Formalmente, \(\forall C \in S\), \(\neg P \in C \implies \Val{\alpha^P}[C] = 1\) perché \(\Val{\alpha^P}[\neg P] = 1\) e 
  \[\neg P \notin C \implies C\setminus\set{P} \in S_P \implies \exists l \in C\setminus\set{P} : \Val{\alpha^P}[l] = \Val{\alpha^P}[l] = 1\]

  Per quanto riguarda \(S_{\neg P}\) il ragionamento è analogo costruendo un assegnamento che assegni \(1\) a \(P\).
  A questo punto possiamo finalmente applicare l'ipotesi induttiva a \(S_P\) e \(S_{\neg P}\), avendo che entrambi sono insoddisfacibili, e quindi \(S_P\vdash_R \square\) e \(S_{\neg P}\vdash_R \square\).

  In particolare quindi esistono due alberi di derivazione del calcolo di risoluzione che hanno uno dei due insiemi come foglie e \(\square\) in radice. Chiamiamoli \(\pi_P\) e \(\pi_{\neg P}\) e valutiamo i diversi casi:
  \begin{itemize}
    \item Se vale \(S_P\subseteq S\) o \(S_{\neg P}\subseteq S\), ovvero una tra \(P\) e \(\neg P\) non è presente in nessuna clausola di \(S\), allora \(S\vdash_R \square\) banalmente per monotònicità.
    \item Alternativamente, esiste almeno una foglia in \(\pi_P\) etichettata con una clausola \(C_P\) e una in \(\pi_{\neg P}\) etichettata con una clausola \(C_{\neg P}\) tali che \(C_P\) e \(C_{\neg P}\) non sono in \(S\).
    In altre parole c'è almeno una foglia per albero deduzione che è stata ottenuta dalla sottrazione del letterale.
    Consideriamo quindi un'applicazione di risoluzione tra \(C_P\) e una qualsiasi clausola \(D \in S_P\) in clash con \(C_P\).
    Chiaramente \(P\notin C_P \land P \notin D \implies P \notin R\) risolvente di \(C_P\) e \(D\).
    A questo punto però aggiungendo \(P\) sia a \(C_P\), che ricordiamo per costruzione originariamente conteneva \(P\) in \(S\), e a \(R\), otterrei ancora un applicazione corretta di risoluzione, poiché applicata a un altro letterale in clash.
    \[\infer[\text{Risoluzione}]{R}{C_P \quad D}\implies\infer[\text{Risoluzione}]{R\cup\set{P}}{C_P\cup\set{P} \quad D}\]
    Con questa idea possiamo aggiungere \(P\) a tutte le foglie di \(\pi_P\) che non appartengono a \(S\), propagando \(P\) ai risolventi delle deduzioni in cui sono coinvolte.
    \begin{center}
      \begin{minipage}{0.42\textwidth}
        \centering
        \begin{prooftree}
          \tiny
          \AxiomC{\(C_P\)}
          \AxiomC{\(D\)}
          \BinaryInfC{\(R\)}
          \AxiomC{\(\vdots\)}
          \BinaryInfC{\(\vdots\)}
          \AxiomC{\(\vdots\)}
          \AxiomC{\(C_P'\)}
          \AxiomC{\(D'\)}
          \BinaryInfC{\(R'\)}
          \BinaryInfC{\(\vdots\)}
          \BinaryInfC{\(\square\)}
        \end{prooftree}
      \end{minipage}
      \hspace{0.3cm}
      \(\xrightarrow{P}\)
      \hspace{0.3cm}
      \begin{minipage}{0.42\textwidth}
        \centering
        \begin{prooftree}
          \tiny
          \AxiomC{\(C_P\cup\set{P}\)}
          \AxiomC{\(D\)}
          \BinaryInfC{\(R\cup\set{P}\)}
          \AxiomC{\(\vdots\)}
          \BinaryInfC{\(\vdots\)}
          \AxiomC{\(\vdots\)}
          \AxiomC{\(C_P'\cup\set{P}\)}
          \AxiomC{\(D'\)}
          \BinaryInfC{\(R'\cup\set{P}\)}
          \BinaryInfC{\(\vdots\)}
          \BinaryInfC{\(\set{P}\)}
        \end{prooftree}
      \end{minipage}
    \end{center}
    In questo modo, abbiamo costruito una nuova deduzione di risoluzione che, a partire da \(S\) essendo tute le foglie etichettate con clausole in \(S\), deduce \(\set{P}\).

    Applicando il ragionamento duale a \(S_{\neg P}\) otteniamo da \(S\) anche una deduzione per \(\set{\neg P}\).
        \begin{center}
      \begin{minipage}{0.42\textwidth}
        \centering
        \begin{prooftree}
          \tiny
          \AxiomC{\(C_{\neg P}\)}
          \AxiomC{\(D\)}
          \BinaryInfC{\(R\)}
          \AxiomC{\(\vdots\)}
          \BinaryInfC{\(\vdots\)}
          \AxiomC{\(\vdots\)}
          \AxiomC{\(C_{\neg P}'\)}
          \AxiomC{\(D'\)}
          \BinaryInfC{\(R'\)}
          \BinaryInfC{\(\vdots\)}
          \BinaryInfC{\(\square\)}
        \end{prooftree}
      \end{minipage}
      \hspace{0.2cm}
      \(\xrightarrow{\neg P}\)
      \hspace{0.2cm}
      \begin{minipage}{0.42\textwidth}
        \centering
        \begin{prooftree}
          \tiny
          \AxiomC{\(C_{\neg P}\cup\set{\neg P}\)}
          \AxiomC{\(D\)}
          \BinaryInfC{\(R\cup\set{\neg P}\)}
          \AxiomC{\(\vdots\)}
          \BinaryInfC{\(\vdots\)}
          \AxiomC{\(\vdots\)}
          \AxiomC{\(C_{\neg P}'\cup\set{\neg P}\)}
          \AxiomC{\(D'\)}
          \BinaryInfC{\(R'\cup\set{\neg P}\)}
          \BinaryInfC{\(\vdots\)}
          \BinaryInfC{\(\set{\neg P}\)}
        \end{prooftree}
      \end{minipage}
    \end{center}
    A questo punto, applicando un ultima volta la regola di risoluzione, è possibile ottenere una deduzione da \(S\) a \(\square\).

    \begin{prooftree}
      \tiny
      \AxiomC{\(C_P\cup\set{P}\)}
      \AxiomC{\(D\)}
      \BinaryInfC{\(R\cup\set{P}\)}
      \AxiomC{\(\vdots\)}
      \BinaryInfC{\(\vdots\)}
      \AxiomC{\(\vdots\)}
      \AxiomC{\(C_P'\cup\set{P}\)}
      \AxiomC{\(D'\)}
      \BinaryInfC{\(R'\cup\set{P}\)}
      \BinaryInfC{\(\vdots\)}
      \BinaryInfC{\(\set{P}\)}
    
      \AxiomC{\(C_{\neg P}\cup\set{\neg P}\)}
      \AxiomC{\(D\)}
      \BinaryInfC{\(R\cup\set{\neg P}\)}
      \AxiomC{\(\vdots\)}
      \BinaryInfC{\(\vdots\)}
      \AxiomC{\(\vdots\)}
      \AxiomC{\(C_{\neg P}'\cup\set{\neg P}\)}
      \AxiomC{\(D'\)}
      \BinaryInfC{\(R'\cup\set{\neg P}\)}
      \BinaryInfC{\(\vdots\)}
      \BinaryInfC{\(\set{\neg P}\)}
      \BinaryInfC{\(\square\)}
    \end{prooftree}
  \end{itemize}
\end{proof}

Per concludere, andiamo a mostrare un esempio della costruzione presente in questa dimostrazione.

\begin{example}{}
  Supponiamo di avere
  \[S = \set{\set{P,Q}, \set{P,\neg Q,\neg R}, \set{R}, \set{\neg P,\neg R}, \set{\neg P,Q,R} }\]
  Fissando \(P\) possiamo creare gli insiemi
  \[S_P=\set{\set{Q},\set{\neg Q, \neg R}, \set{R}}\]
  \[S_{\neg P} = \set{\set{R}, \set{\neg R}, \set{Q,R}}\]

  Ora è facile vedere che \(S_P \models_R\square\) e \(S_{\neg P} \models_R \square\):
  \begin{prooftree}
    \tiny
    \AxiomC{\(\set{Q}\)}
    \AxiomC{\(\set{\neg Q, \neg R}\)}
    \BinaryInfC{\(\set{\neg R}\)}
    \AxiomC{\(\set{R}\)}
    \BinaryInfC{\(\square\)}
  \end{prooftree}

  \begin{prooftree}
    \tiny
    \AxiomC{\(\set{R}\)}
    \AxiomC{\(\set{\neg R}\)}
    \BinaryInfC{\(\square\)}
  \end{prooftree}

  Per queste deduzioni abbiamo che \(\set{\neg R}\) in \(S_{\neg P}\) non apparteneva a \(S\), così come \(\set{\neg Q, \neg R}\).
  Aggiungiamo quindi alle deduzioni \(P\) e \(\neg P\), e propaghiamole.

  \begin{prooftree}
    \tiny
    \AxiomC{\(\set{Q}\)}
    \AxiomC{\(\set{\neg Q, \neg R,P}\)}
    \BinaryInfC{\(\set{\neg R,P}\)}
    \AxiomC{\(\set{R}\)}
    \BinaryInfC{\(\set{P}\)}
  \end{prooftree}

  \begin{prooftree}
  \tiny
    \AxiomC{\(\set{R}\)}
    \AxiomC{\(\set{\neg R,\neg P}\)}
    \BinaryInfC{\(\set{\neg P}\)}
  \end{prooftree}

  A questo punto possiamo dedurre \(\square\) facilmente col clash di \(P\) e \(\neg P\).

  \begin{prooftree}
    \tiny
    \AxiomC{\(\set{Q}\)}
    \AxiomC{\(\set{\neg Q, \neg R,P}\)}
    \BinaryInfC{\(\set{\neg R,P}\)}
    \AxiomC{\(\set{R}\)}
    \BinaryInfC{\(\set{P}\)}
    \AxiomC{\(\set{R}\)}
    \AxiomC{\(\set{\neg R,\neg P}\)}
    \BinaryInfC{\(\set{\neg P}\)}
    \BinaryInfC{\(\square\)}
  \end{prooftree}
\end{example}